% !TEX program = xelatex
\documentclass[11pt,a4paper]{ctexart}

\usepackage[margin=2.3cm]{geometry}
\usepackage{amsmath,amssymb}
\usepackage{booktabs}
\usepackage{longtable}
\usepackage{array}
\usepackage{multirow}
\usepackage{enumitem}
\usepackage[hidelinks]{hyperref}
\usepackage{xcolor}
\usepackage{caption}
\hypersetup{colorlinks=true, linkcolor=black, urlcolor=blue, citecolor=black}
\captionsetup{font=small,labelfont=bf}
\setlength{\parskip}{0.35em}
\setlength{\parindent}{2em}
\setlength{\emergencystretch}{3em}
\setlength{\tabcolsep}{4pt}
\setlist{nosep}
\renewcommand{\arraystretch}{1.15}

\newcommand{\confhigh}{\textbf{高}}
\newcommand{\confmid}{中}
\newcommand{\conflow}{\textit{低}}

\title{\bfseries 趋势跟踪与均值回归交易规则设计\\[3pt]
\large 基于《趋势跟踪与均值回归：同一自相关结构的两侧剪影》的实证结论}
\author{叨客宁宁}
\date{}

\begin{document}

\maketitle

\begin{abstract}
\noindent 本文把综述《趋势跟踪与均值回归：同一自相关结构的两侧剪影》的实证结论翻译为可执行的交易规则，
共三套：\textbf{期货时序趋势跟踪}（对应论文中期货时序口径趋势全窗口占优的证据）、
\textbf{A 股截面反转}（对应 A 股个股截面反转在 $L=20\sim120$ 全面占优、$L=60$ 净夏普 $0.99$ 的证据）、
\textbf{期货截面统计套利}（对应期货截面反转在 $L=60/120$ 占优、净夏普 $0.72/0.56$ 且
Newey--West $t=2.59$ 的证据）。

三套规则共享同一套数据、成本与可交易性口径，均按"适用性依据 $\to$ 信号 $\to$ 组合构建 $\to$
仓位与缩放 $\to$ 进出场纪律 $\to$ 止损与失效防线 $\to$ 成本门槛 $\to$ 容量约束"的固定结构展开，
每条规则附\textbf{参数边界表}（说明每个参数在论文实证中的取值范围与越界后果）
与\textbf{证据映射表}（逐条对应论文结论、证据来源与置信度）。

三条贯穿全部规则的设计铁律：\textbf{其一，每条规则必须能追溯到论文的具体证据}，
凡属工程改进而非本文实证支持的做法，一律显式标注；
\textbf{其二，成本先于信号}——任何信号的预期毛收益必须超过其往返成本的 3 倍，
否则该参数组合直接废弃（论文实证中 $L=5$ 的 A 股反转毛利夏普 $1.08$ 被 $12.31\%$ 的成本拖累压至 $-0.35$，是这条铁律的来源）；
\textbf{其三，失效防线内置}——趋势规则内置恐慌状态过滤与拥挤监控，反转规则内置止损与资金冗余，
因为两策略的失效模式不同（前者"慢性出血加跳崖"，后者"不收敛加资金断裂"）。

本文不宣称任何规则能产生正收益。论文已明确其自建实证是\emph{变量受控的相对比较}而非绝对收益估计；
本文沿用了这一限定，所有规则的可执行性以论文披露的边界为限。

为把「可执行」落到实处，本文另附三份东西。第一，\textbf{可运行的回测代码框架}
（\texttt{code/} 目录）：规则一、规则二各一套完整回测实现，共享同一数据、成本与评估口径；
代码已实跑，三处实现缺陷（清仓止损的永久锁定、对冲 beta 未按因果估计、Kalman 使用当期观测
制造虚假均值回归）在实跑中暴露并修正，全部过程如实记录。
第二，\textbf{规则三的协整配对筛选算法}（第~\ref{sec:pairalgo}~节）：
七步流程、每步判据与实测拒绝率；在中国期货 16992 次配对检验中仅 75 条通过（通过率 $0.44\%$），
印证了反转策略「条件性」的结构特征（改用 Engle--Granger 专用临界值后为 26 条，见第~\ref{sec:verify}~节）。
第三，\textbf{上线前的多重检验校正流程}（第~\ref{sec:mtflow}~节）：
四道闸门（White Reality Check、Hansen SPA、紧缩夏普比率、过拟合概率）对 60 个参数组合实跑，
结果出现分歧——自举类检验边际拒绝，而紧缩夏普比率判定最优组合的幅度不足以排除数据窥探。
本文如实保留这一分歧并给出解读。

第四，\textbf{四项验证补充}（第~\ref{sec:verify}~节）：
样本外 walk-forward（样本内夏普 $1.26$ 衰减到样本外 $0.57$，仅 $44\%$ 的折为正）、
行业中性化口径对照（真实申万 vs 聚类代理，量级一致）、
规则一完整窗口网格（推翻了前文对 $L=250$ 的推荐）、
规则三统计口径修正（Engle--Granger 专用临界值使通过数减少 $65\%$，严格口径下成交仅 $7$ 笔）。
这四项验证各自修正了前文的一处判断，共同把本文的定位收紧为
\textbf{设计模板与验证流程}，而非可上线策略。

第五，\textbf{依据上述验证重标定的参数}（第~\ref{sec:retune}~节）：
规则一的窗口由「$250/60/10$ 三档」改为 $\mathbf{L=10}$（首选）、$\mathbf{L=20}$（备选），
且窗口必须与入场门槛\emph{联合}选择（$L=250$ 在 $K=1.5$ 处的孤立峰按稳健性判据拒收）；
规则二的窗口由 $L=60$ 改为 $\mathbf{L=250}$（首选）、$L=60$（备选），
依据是二者夏普几乎相同而 $L=250$ 的最大回撤浅 $16$ 个百分点、成本占毛利低 $6.6$ 个百分点；
ATR 止损倍数被实测判定为\emph{低敏感参数}（夏普跨度 $0.04$），无须精调。
必须同时声明的是：重标定改善的是\emph{相对排序}，
两条规则的首选参数在自身 walk-forward 中的样本外夏普分别为 $-0.03$ 与 $-0.18$，\textbf{均不兑现正收益}。
\end{abstract}

\tableofcontents
\newpage

\section{引言：规则从哪里来}\label{sec:intro}

\subsection{问题}

上一份综述得出的核心结论是：趋势跟踪与均值回归共享同一个信号，唯一区别是符号；两者的分野不在信号里，而在信号之后的每一步。这个结论对“设计规则”这件事提出了一个苛刻的要求——如果两策略的信号相同，那么规则设计的重心就\emph{不应该}放在“找一个更好的信号”上，而应该放在验证、缩放、对冲、交易与融资这五个环节上。

本文据此设计三套规则。选择这三套而不是别的，依据是论文实证中“哪几个口径曾经给出过正收益”：

\begin{itemize}[leftmargin=2.2em]
  \item \textbf{期货时序趋势}：论文表 5 显示，在中国期货时序口径下，趋势策略在 $L=5$ 到 $L=250$ 的全部六个窗口上净夏普为正（$0.09\sim0.32$），而时序反转全为负（$-0.17\sim-0.60$）。这是论文中趋势一侧唯一全窗口成立的口径。
  \item \textbf{A 股截面反转}：论文表 6 显示，A 股个股截面反转在 $L=20$ 到 $L=250$ 上净夏普正（$0.56\sim0.99$），其中 $L=60$ 净夏普 $0.99$、年化 $8.98\%$、最大回撤 $-18.64\%$、月胜率 $60\%$。这是论文全部实验中夏普最高的一组。
  \item \textbf{期货截面统计套利}：论文表 5 显示，期货截面反转在 $L=60$ 与 $L=120$ 上净夏普 $0.72$ 与 $0.56$，且 $L=60$ 的 Newey--West $t$ 统计量为 $2.59$——这是论文中少数通过常规显著性门槛的结果。
\end{itemize}

其余口径（A 股宽基指数时序）论文已明确标注证据力弱，本文不为它设计规则。

\subsection{三条设计铁律}

\subsubsection*{铁律一：每条规则必须能追溯到论文的具体证据}

交易规则最常见的失效方式不是逻辑错误，而是\emph{无据可依}——某个参数取 $60$ 日、某个阈值取 $2$ 倍标准差，说不出为什么。本文要求每套规则中的每个关键决策都指回论文中的具体表格或结论。

这条铁律带来一个副作用，必须提前说明：\textbf{论文没有检验过的东西，本文不得声称为有效}。例如多窗口信号集成、ATR 止损的具体倍数、参与率约束的冲击成本模型，这些在论文中或未检验、或仅作为执行章的定性讨论出现。本文会给出这些做法的工程建议，但会在证据映射表中把它们标为\emph{工程判断}而非\emph{实证支持}，与有论文支撑的条目明确分开。

\subsubsection*{铁律二：成本先于信号}

论文实证中最刺眼的一条证据是 A 股截面反转在 $L=5$ 上的表现：毛利夏普 $+1.08$，净夏普 $-0.35$。信号本身相当强（毛夏普超过 1），但 $79.8$ 倍的年化换手带来 $12.31\%$ 的年化成本拖累，把信号完全吞掉。

这条证据确立了本文的成本门槛：

\begin{quote}
\itshape
任何参数组合，只有当预期毛收益超过其往返成本的 3 倍时，才允许进入考察范围。
达不到这条线的组合，无论其信号统计量多显著，一律废弃。
\end{quote}

门槛设为 3 倍而非更高，是因为论文中表现最好的一档（期货截面反转 $L=60$）的成本占毛利比约为 $6.3\%$，即毛利约为成本的 16 倍；而 $L=60$ 的 A 股反转成本占毛利 $31.87\%$，约 3.1 倍——恰好在门槛线上。换言之，3 倍这条线是“论文实证中最勉强可用的组合”所对应的位置。它不是一个理论最优值，而是一个从实证中反推出来的下限。

\subsubsection*{铁律三：失效防线内置}

论文的对偶六指出两策略的失效模式不同：趋势策略是“震荡市慢性出血叠加崩溃时跳崖”，反转策略是“不收敛加资金断裂、跳变式亏损”。防线因此不能通用。

\begin{itemize}[leftmargin=2.2em]
  \item 趋势规则内置两道防线：\textbf{恐慌状态过滤}（对移动量崩溃）与\textbf{拥挤监控}（对信号同质化）。
  \item 反转规则内置两道防线：\textbf{偏离止损}（对不收敛）与\textbf{资金冗余}（对被迫平仓）。
\end{itemize}

需要特别强调一条反直觉的纪律：\textbf{反转策略不能用“越跌越买”的逻辑加仓}。论文在处理 LTCM 与 2007 年量化危机时指出，反转策略的亏损本身不可怕，可怕的是亏损触发了资金约束——业绩与资金流的正反馈会让套利者在错价最深时被迫减仓\cite{shleifervishny1997,khandanilo2011}。因此反转规则的加仓必须绑定在“价差已开始收敛”的证据上，而不是绑在“偏离更大了”上。

\subsection{三套规则的共同骨架}

本文对每套规则采用同一套叙述结构，便于横向比较与逐项复核：

\begin{enumerate}[leftmargin=2.2em]
  \item \textbf{适用性依据}：论文中支持该规则的具体证据；
  \item \textbf{信号}：计算什么、用什么窗口、如何处理噪声；
  \item \textbf{组合构建}：选多少标的、如何加权、如何中性化；
  \item \textbf{仓位与缩放}：波动率目标化是否适用（三套规则答案不同）；
  \item \textbf{进出场纪律}：调仓频率、免交易带、入场门槛；
  \item \textbf{止损与失效防线}：针对该策略失效模式的防线；
  \item \textbf{成本门槛与容量}：成本核算、容量约束来源；
  \item \textbf{参数边界表}：每个参数的取值、论文依据、越界后果；
  \item \textbf{证据映射表}：决策 $\to$ 论文结论 $\to$ 证据来源 $\to$ 置信度。
\end{enumerate}

第~\ref{sec:shared}~节先给出三套规则共享的基础设施（数据口径、成本模型、可交易性过滤、评估指标），此后三节各展开一套规则。

\subsection{边界声明}

本文是规则设计文档，不是策略推介。三条限定必须写在前面。

\textbf{其一，绝对收益不可预期}。论文的自建回测是变量受控的相对比较，未做多重检验校正（$6$ 个窗口 $\times$ 两个方向 $\times$ 四组口径 $=48$ 个组合），且 A 股侧做空不可实现、除权修复存在残余污染。本文沿用的所有绩效数字（如 $L=60$ 的 $8.98\%$ 年化）都应读作“在该口径下的相对位置”，不是可实现的收益预期。

\textbf{其二，A 股反转的多空版本不可直接执行}。论文披露 A 股融券受限、成本高、可借券池小，截面多空组合在多数时段不可实现。本文因此为 A 股反转给出两个版本：可执行的“纯多头加指数对冲”版本，与研究用的“多空”版本，后者明确标注不可实盘使用。

\textbf{其三，制度环境不可外推}。论文样本期为 2014--2026 年，其间中国期货与 A 股经历了股指期货交易限制、融券政策调整、注册制改革等制度变化。规则的适用性以该期间为限。

\section{共同基础设施}\label{sec:shared}

三套规则共享同一套数据与成本口径。这不是形式上的统一，而是论文结论本身的要求：论文所有比较的合法性都建立在“同一资产池、同一调仓频率、同一成本模型、同一评估指标、唯一区别是信号符号”这一设定上。若三套规则各自使用不同的成本假设，本文的规则之间就无法与论文的数字对照。

\subsection{数据口径}

\textbf{期货}。使用主力连续合约，并经换月跳空守门处理：先检测候选跳空（日收益超过前移滚动标准差的 10 倍且绝对值超过 2\%，且落在主力切换日 $\pm2$ 交易日窗口内），再用新主力合约当日收益或加权指数交叉确认，确认后做前复权调整。论文实证中的期货侧数据即按此处理。

\textbf{A 股}。通达信本地日线为不复权价，论文的修复规则是：主板个股单日收益低于 $-10.5\%$、创业板与科创板低于 $-20.5\%$ 者判定为除权，将该日收益置零（等价于按当日比例对更早价格做前复权）。论文样本内共识别并修复 $6121$ 处。

\textbf{必须补齐的缺口}。论文在审计章明确披露这项修复的缺陷：\emph{跳空幅度小于涨跌停幅度的除权无法识别}（如小额分红、$10$ 送 $1$），构成残余污染。本文的规则若用于实盘，必须补上这一层——可行做法是从数据商处取得复权因子序列，用它替换论文的近似规则。在补上之前，A 股反转规则的所有历史回测数字都带有这一已知偏差。

\textbf{指数}。A 股宽基指数仅用于两处：作为反转组合的对冲工具与业绩基准。论文已指出 7 个 A 股宽基指数的相关系数高达 $0.85$ 以上，其“组合”实质是单一 beta，不能用于分散。

\subsection{成本模型}

\begin{table}[htbp]
\centering\footnotesize
\caption{成本参数（与论文一致）}\label{tab:cost}
\begin{tabular}{@{}>{\raggedright\arraybackslash}p{4.6cm}>{\raggedright\arraybackslash}p{2.2cm}>{\raggedright\arraybackslash}p{7.5cm}@{}}
\toprule
\textbf{市场} & \textbf{单边成本} & \textbf{口径说明} \\
\midrule
期货 & 2 个基点 & 含手续费与滑点；论文口径 \\
A 股个股 & 15 个基点 & 含佣金、印花税（卖出千分之一）与滑点 \\
A 股指数/ETF & 5 个基点 & 对应 ETF 或股指期货实现 \\
\bottomrule
\end{tabular}
\end{table}

\textbf{成本计算规则}（沿用论文）：只对换入与换出部分收费，持有未动的仓位零成本。这在权重随收益漂移的设定下意味着：即使目标权重不变，实际权重已偏离，因此每次调仓仍会产生小额换手。

\textbf{大额资金必须升级成本模型}。论文的固定基点模型对反转类规则\emph{系统性低估}成本，原因是反转交易密集发生在流动性较差的时点（偏离越大时往往流动性越差，即流动性螺旋\cite{brunnermeier2009}）。实盘应改用参与率约束的冲击模型：

\begin{equation}
\text{冲击成本} = K\cdot\sqrt{\text{参与率}},\qquad
\text{参与率}=\frac{\text{本笔委托量}}{\text{同期市场成交量}} ,
\end{equation}

并设定参与率上限（论文工程约定为 $10\%$），超过上限的委托必须拆分或标记拒绝。$K$ 需按品种用历史成交数据标定。这一升级在论文中属定性讨论，本文将其列为\textbf{工程判断}。

\subsection{可交易性过滤}

三套规则共用同一套过滤，避免回测高估：

\begin{itemize}[leftmargin=2.2em]
  \item \textbf{涨跌停}：主板 $\pm10\%$，创业板（$300$/$301$）、科创板（$688$/$689$）$\pm20\%$。涨跌停当日无法成交，该标的的调仓顺延至下一可交易日。
  \item \textbf{停牌}：停牌期间不调仓；复牌首日不交易（价格发现未完成）。
  \item \textbf{流动性门槛}：期货要求近 20 日日均成交额不低于 1 亿元；A 股个股要求不低于 5000 万元。低于门槛者从候选池剔除。
  \item \textbf{次新与特殊处理}：A 股剔除 ST 与 *ST，剔除上市不足 250 个交易日的次新股。
\end{itemize}

论文披露其回测\emph{未做}涨跌停与停牌过滤，因此在高波动时段（如 2015 年）会高估可实现收益。本文把这些过滤写成规则的一部分，属于对论文局限的修补，标注为\textbf{工程判断}。

\subsection{评估指标}

三套规则统一报告 13 项指标，与论文一致：年化收益、年化波动、Sharpe、最大回撤、Calmar、偏度、超额峰度、年化换手、毛利年化、成本拖累年化、月胜率、Newey--West $t$ 统计量（5 阶滞后）、成本占毛利比。

其中三项在使用时需特别注意：

\textbf{成本占毛利比}是本文最看重的指标，因为它直接对应铁律二。论文中该比值从 $L=250$ 期货趋势的 $10.2\%$ 到 $L=5$ 期货趋势的 $50.1\%$ 不等；超过 $33\%$（即毛利不足成本 3 倍）的组合按铁律二废弃。

\textbf{偏度与超额峰度}用于判断策略的尾部特征。论文指出趋势策略总体右偏但含深度左尾（动量崩溃），反转策略左偏。规则上线前应检查实测偏度是否与预期方向一致；若趋势规则的实测偏度显著为负，说明它已退化为“没有凸性保护的追涨”。

\textbf{Newey--West $t$ 统计量}用于初步显著性判断。注意论文未做多重检验校正，因此 $t>2$ 只能视为“值得进一步检验”，不能视为“已证实显著”。论文中期货截面反转 $L=60$ 的 $t=2.59$ 是少数超过 2 的结果。

\section{规则一：期货时序趋势跟踪}\label{sec:trend}

\subsection{适用性依据}

论文表 5 的期货时序口径（TS）结果：

\begin{table}[htbp]
\centering\footnotesize
\caption{期货时序口径下趋势策略的实测（论文表 5，成本 2 基点）}\label{tab:trendsrc}
\begin{tabular}{@{}>{\raggedright\arraybackslash}p{1.2cm}>{\raggedright\arraybackslash}p{1.8cm}>{\raggedright\arraybackslash}p{1.8cm}>{\raggedright\arraybackslash}p{1.8cm}>{\raggedright\arraybackslash}p{1.8cm}>{\raggedright\arraybackslash}p{1.8cm}>{\raggedright\arraybackslash}p{1.8cm}@{}}
\toprule
\textbf{$L$} & \textbf{净夏普} & \textbf{毛利夏普} & \textbf{年化收益} & \textbf{最大回撤} & \textbf{年化换手} & \textbf{成本占毛利} \\
\midrule
5   & $+0.20$ & $+0.40$ & $+0.92\%$ & $-13.97\%$ & $52.41$ & $50.1\%$ \\
10  & $+0.32$ & $+0.45$ & $+1.56\%$ & $-14.23\%$ & $35.14$ & $29.4\%$ \\
20  & $+0.21$ & $+0.30$ & $+1.01\%$ & $-15.24\%$ & $24.78$ & $30.0\%$ \\
60  & $+0.21$ & $+0.26$ & $+1.06\%$ & $-14.96\%$ & $16.14$ & $20.9\%$ \\
120 & $+0.09$ & $+0.13$ & $+0.36\%$ & $-22.78\%$ & $12.10$ & $31.0\%$ \\
250 & $+0.24$ & $+0.27$ & $+1.36\%$ & $-15.97\%$ & $8.83$ & $10.2\%$ \\
\bottomrule
\end{tabular}
\end{table}

三个可直接读出的判断。其一，六个月窗口全部为正收益，趋势侧在这个口径上稳健。其二，$L=10$ 夏普最高（$0.32$），但换手 $35.14$ 倍、成本占毛利 $29.4\%$，已接近铁律二的 $33\%$ 边界。其三，$L=250$ 是成本效率最好的一档——换手仅 $8.83$ 倍、成本占毛利 $10.2\%$，夏普 $0.24$ 且年化 $1.36\%$ 为全窗口最高。这构成“低换手档”的实证依据。

需要如实指出一点：这一档的绝对收益很低（年化 $1.36\%$）。论文的期货趋势结果整体弱于文献中的国际样本（Hurst 等 2017 的百年样本、Moskowitz 等 2012 的 58 合约）\cite{mop2012,hop2017}，这与样本期仅 2014--2026 年、且中国商品期货该期间趋势性偏弱有关。规则设计必须以这个较弱的基准为出发点，不能引用国际文献的收益量级。

另外，本规则的信号构造（过去 $L$ 期累计收益的符号、按波动率缩放、跨品种组合）直接沿用时间序列动量的标准范式\cite{mop2012}，并在百年尺度上被验证过稳健性\cite{hop2017}。这意味着规则的\emph{形式}有充分的文献支持，但\emph{在中国期货上的收益水平}只由论文表 5 支撑——两者不能混为一谈。

\subsection{信号}

\textbf{基础信号}。对品种 $i$ 在交易日 $t$：

\begin{equation}
s_i(L)=\sum_{k=0}^{L-1}\ln\left(1+r_{i,t-k}\right),\qquad
\sigma_i(L)=\mathrm{sign}\bigl(s_i(L)\bigr).
\end{equation}

\textbf{窗口选择}。论文按成本效率给出三档（$L=10$ 高效、$L=60$ 平衡、$L=250$ 低换手），但那是在没有入场门槛与止损的口径下读出的。加入本规则的工程增强项后重做网格，排序发生实质变化（重标定依据见第~\ref{sec:retune}~节），本文因此给出两档：

\begin{itemize}[leftmargin=2.2em]
  \item \textbf{首选 $L=10$}：工程增强后夏普 $0.44$、成本占毛利 $18.1\%$，参数曲面随门槛单调平滑，稳健性最好；
  \item \textbf{备选 $L=20$}：夏普 $0.33$、最大回撤 $-16.9\%$，适合更看重回撤的资金。
\end{itemize}

$L=250$ 与 $L=120$ 在工程增强下失效（第~\ref{sec:verify}~节：$L=250$ 的夏普由 $0.24$ 落到 $-0.00$、成本占毛利升至 $108.7\%$），$L=5$ 在 $K\ge1.0$ 时不通过成本门槛。窗口还必须与入场门槛\emph{联合}选择——$L=250$ 在 $K=1.5$ 处的孤立峰按稳健性判据拒收。

\textbf{入场门槛}。论文指出趋势信号的噪声在短窗口上极大。为过滤无意义的微小信号，要求标准化信号强度超过噪声水平才允许建仓：

\begin{equation}
\frac{\left|s_i(L)\right|}{\hat\sigma_i\sqrt{L}}>\;1.0 ,
\end{equation}

其中 $\hat\sigma_i$ 为日收益标准差。不满足门槛的品种当日目标权重置零。这一门槛在论文中未检验，属\textbf{工程判断}。

\textbf{关于多窗口集成}。工程上有一种常见做法是把多个窗口的符号加权平均（如 $L\in\{10,20,60,120,250\}$ 等权投票），以期降低单一窗口的参数风险。论文\emph{未}检验集成效果，因此本文不建议把它作为默认做法，而是列为可选增强项，并要求上线前单独回测验证。这一点必须说清楚：集成听起来稳妥，但它属于“未经本文实证支持的做法”。

\subsection{组合构建}

\textbf{候选池}。取流动性分层后日均成交额前 $N$ 的品种，$N$ 建议 $20\sim30$。论文实证使用了 45 个品种的全池，未做品种筛选；剔除低流动性品种会改变组合的风险特征，需单独验证。

\textbf{权重}。时序口径下每个品种独立建仓，权重由波动率与方向共同决定（见下节），归一化使总敞口为 1。

\textbf{为何不做行业中性化}。期货的“行业”即产业链，趋势策略在产业链内的同向暴露是其收益来源之一（即论文所述“持续单向运动”环境下的板块共振），因此不对期货趋势做板块中性化。这一点与 A 股反转规则相反——那里必须做中性化（详见第~\ref{sec:ashare}~节）。差异的根据是论文对偶四：趋势的收益来自承担共同方向风险，反转的收益来自剥离共同因子后的特质回复。

\subsection{仓位与缩放}

\textbf{波动率目标化适用}。这是三套规则中唯一适用波动率目标化的一套。论文引用 Barroso 与 Santa-Clara(2015) 把动量夏普从 $0.53$ 提到 $0.97$、Harvey 等(2018) 在 60 多个资产上验证波动率目标改善风险资产夏普\cite{barroso2015,harvey2018}，并引 Moreira 与 Muir(2017) 的波动率管理组合\cite{moreiramuir2017}。

权重公式：

\begin{equation}
w_i=\sigma_i(L)\cdot\min\left(\frac{V_{\text{target}}}{\hat\sigma_i^{(60)}},\;C_{\text{cap}}\right),
\qquad\text{随后归一化}\quad w\leftarrow\frac{w}{\sum_j|w_j|},
\end{equation}

其中 $\hat\sigma_i^{(60)}$ 为 60 日已实现波动年化值，$V_{\text{target}}=10\%$，$C_{\text{cap}}=3.0$（单品种杠杆上限）。这三个参数取自论文实证设定，与其期货时序回测一致，因此规则的历史绩效可直接与论文表 5 对照。

\textbf{归因警告}。论文引用 Kim 等(2016) 指出，时间序列动量策略的收益有很大一部分来自波动率缩放本身而非时序符号\cite{kim2016}，Huang 等(2020) 更直接地指出 TsMOM 的表现与“按历史样本均值投资”的简单规则几乎相同\cite{huang2020}。这对规则设计有一个直接要求：\textbf{上线前必须做缩放与信号的分离归因}——分别测试“只保留符号、不做缩放”与“只做缩放、符号恒定”两个对照组，确认收益来源。若收益主要来自缩放，则该规则的经济含义应重新表述，而不是继续称为“趋势跟踪”。

\subsection{进出场纪律}

\textbf{调仓频率}：每 5 个交易日。信号使用至 $t$ 日收盘的信息，仓位自 $t+1$ 日起生效。

\textbf{免交易带}（no-trade band）。论文实证显示，即使目标权重不变，持仓随收益漂移仍会产生换手。为压低成本，设免交易带：

\begin{equation}
\text{若}\;\left|w_i^{\text{new}}-w_i^{\text{current}}\right|<\theta_{\text{band}},
\quad\text{则}\;w_i^{\text{new}}\leftarrow w_i^{\text{current}},
\qquad\theta_{\text{band}}=0.1\times\frac{V_{\text{target}}}{\hat\sigma_i^{(60)}} .
\end{equation}

带宽随波动率自适应：波动率高的品种带宽更窄（更频繁调整以免敞口失控），波动率低的品种带宽更宽（减少无谓交易）。这一设计在论文中未检验，属\textbf{工程判断}，但其必要性由论文的成本证据支持——成本占毛利比在短窗口上已达 $29.4\%$，任何能压低换手的机制都有价值。

\subsection{失效防线}

\subsubsection*{防线一：恐慌状态过滤（对移动量崩溃）}

论文引用 Daniel 与 Moskowitz(2016)：动量崩溃发生在恐慌状态（市场下跌后、波动率高企并伴随反弹），且部分可预测；据此构造的动态动量策略把 alpha 与夏普大约翻倍\cite{danielmoskowitz2016}。

\textbf{触发定义}：

\begin{equation}
\text{恐慌状态}\iff
\Bigl(\underbrace{R_{\text{mkt}}^{(60)}<-5\%}_{\text{市场 60 日下跌}}\Bigr)
\;\wedge\;
\Bigl(\underbrace{\sigma_{\text{mkt}}^{(60)}>\text{历史 80 分位}}_{\text{波动率高企}}\Bigr) ,
\end{equation}

其中 $R_{\text{mkt}}^{(60)}$ 为等权市场组合的 60 日收益，$\sigma_{\text{mkt}}^{(60)}$ 为其 60 日已实现波动年化值，分位数在滚动 3 年窗口上计算。

\textbf{动作}：触发后总敞口乘以 $0.5$，持续至波动率回落至 60 分位以下。

\textbf{参数边界}：$-5\%$ 与 80 分位是两个可调参数。设得过严（如 $-10\%$ 且 90 分位）会导致几乎从不触发，防线形同虚设；设得过松（如 $0\%$ 且 60 分位）会导致频繁降仓、错失趋势收益。论文仅给出“恐慌状态”的定性刻画，未给出精确阈值，因此这两个数字属\textbf{工程判断}，需用滚动回测标定。

\subsubsection*{防线二：拥挤监控（对信号同质化）}

论文对偶五指出趋势策略的容量约束来自\emph{拥挤}而非成本：Stein(2009) 的协调失败理论\cite{stein2009}、Lou 与 Polk(2022) 的 comomentum 度量（高拥挤预示动量收益下降）\cite{loupolk2022}、Baltas(2019) 的另类风险溢价拥挤研究\cite{baltas2019}。

\textbf{指标}：策略内品种日收益的两两平均相关系数 $\bar\rho_t$，在滚动 3 年窗口上取分位。

\textbf{动作}：$\bar\rho_t$ 超过 90 分位时，总敞口乘以 $0.7$，并输出预警。

\textbf{边界}：论文明确指出拥挤度的定量估计尚无统一口径，容量上限的估计分歧较大（见论文遗留课题 Q3）。因此本防线只作为\emph{减速器}而非\emph{开关}使用——它降低敞口，但不强制清仓。把它当作精确的容量计量工具是误用。

\subsubsection*{防线三：组合止损}

\begin{itemize}[leftmargin=2.2em]
  \item 组合最大回撤达 $15\%$：总敞口腰斩；
  \item 组合最大回撤达 $25\%$：全部平仓，暂停 20 个交易日后再评估；
  \item 单品种浮亏达 $2.5\times\text{ATR}_{20}$：该品种平仓。
\end{itemize}

论文实验中的最大回撤为 $-13.97\%$ 至 $-22.78\%$，因此 $15\%$ 与 $25\%$ 两条线分别位于历史回撤的浅部与深部，属于\emph{基于实测回撤分布}的设定。ATR 倍数属\textbf{工程判断}。

\textbf{不得使用反转式加仓}。趋势策略在震荡市中的常规应对是止损离场而非逢低加仓；论文指出趋势策略的失效包含“慢性出血”，加仓会把慢性出血放大为急性亏损。

\subsection{成本门槛与容量}

\textbf{成本核算}。按表~\ref{tab:cost}，期货单边 2 基点。按铁律二检查：

\begin{table}[htbp]
\centering\footnotesize
\caption{期货趋势各窗口的成本门槛检查（论文口径）}\label{tab:trendcost}
\begin{tabular}{@{}>{\raggedright\arraybackslash}p{1.6cm}>{\raggedright\arraybackslash}p{2.6cm}>{\raggedright\arraybackslash}p{2.6cm}>{\raggedright\arraybackslash}p{3.6cm}@{}}
\toprule
\textbf{$L$} & \textbf{成本占毛利} & \textbf{毛利/成本} & \textbf{门槛判定（要求 $>3$）} \\
\midrule
250 & $10.2\%$ & $9.8$ & 通过（余量充足） \\
60  & $20.9\%$ & $4.8$ & 通过 \\
120 & $31.0\%$ & $3.2$ & 勉强通过 \\
20  & $30.0\%$ & $3.3$ & 勉强通过 \\
10  & $29.4\%$ & $3.4$ & 勉强通过 \\
5   & $50.1\%$ & $2.0$ & \textbf{不通过，废弃} \\
\bottomrule
\end{tabular}
\end{table}

这张表给出一个明确的操作结论：\textbf{在 2 基点成本下，$L=5$ 必须废弃}；$L=10/20/120$ 处于门槛边缘，一旦实际成本升到 3 基点以上（例如加入冲击成本），应优先切换到 $L=60$。需要提醒的是，本表按\emph{论文口径}计算；本规则的工程增强会改变各窗口的成本占毛利（实测值见第~\ref{sec:verify}~节的表~\ref{tab:grid1}），其中 $L=250$ 的余量优势在增强后消失、成本占毛利反升至 $108.7\%$，因此“低换手档”这一档在增强口径下不再可用。

\textbf{容量约束}。趋势的容量约束来自拥挤，不来自成本。论文未给出容量上限的定量估计，本文因此不给数字，只给监控项：comomentum 分位、策略资金规模占品种持仓量的比例、以及与同类 CTA 的相关性。

\subsection{参数边界表}

\begin{table}[htbp]
\centering\footnotesize
\caption{规则一参数边界}\label{tab:trendparam}
\begin{tabular}{@{}>{\raggedright\arraybackslash}p{2.5cm}>{\raggedright\arraybackslash}p{2.2cm}>{\raggedright\arraybackslash}p{3.3cm}>{\raggedright\arraybackslash}p{6.4cm}@{}}
\toprule
\textbf{参数} & \textbf{取值} & \textbf{证据来源} & \textbf{越界后果} \\
\midrule
窗口 $L$ & $10$（首选）/ $20$（备选） & 论文表 5 + 第~\ref{sec:retune}~节 & 短于 $10$ 成本占毛利超 $33\%$ 须废弃；$250$ 与 $120$ 在工程增强下失效 \\
调仓频率 & 5 交易日 & 论文设定 & 缩短至日频换手翻倍，成本占毛利可能翻倍 \\
目标波动率 & 10\% & 论文设定（MOP12/HOP17 体系） & 提高则杠杆上升、回撤放大 \\
单品种杠杆上限 & 3.0 & 论文设定 & 提高则单品种主导组合风险 \\
免交易带系数 & 0.1 & 工程判断 & 加宽则跟踪滞后；收窄则换手上升 \\
入场门槛 & 1.0 & 工程判断 & 提高则持仓数骤减、集中度上升 \\
恐慌阈值 & $-5\%$ / 80 分位 & DM16 定性 + 工程标定 & 过严则永不触发；过松则错失趋势 \\
拥挤阈值 & 90 分位 / 乘 0.7 & Lou--Polk、Baltas 定性 & 论文注明容量定量估计无统一口径，不宜作精确开关 \\
回撤止损线 & $15\%$ / $25\%$ & 论文实测回撤分布 & 过紧则频繁触发；过松则失去保护 \\
\bottomrule
\end{tabular}
\end{table}

\subsection{证据映射表}

\begin{table}[htbp]
\centering\footnotesize
\caption{规则一：决策与论文证据的映射}\label{tab:trendmap}
\begin{tabular}{@{}>{\raggedright\arraybackslash}p{3.4cm}>{\raggedright\arraybackslash}p{5.3cm}>{\raggedright\arraybackslash}p{2.6cm}>{\raggedright\arraybackslash}p{2.6cm}@{}}
\toprule
\textbf{决策} & \textbf{论文结论} & \textbf{证据来源} & \textbf{置信度} \\
\midrule
用趋势方向而非反转 & 期货时序趋势全窗口夏普为正、反转全为负 & 论文表 5 & \confhigh \\
不用 $L=5$ & $L=5$ 成本占毛利 $50.1\%$ & 论文表 5 & \confhigh \\
波动率目标化 & 趋势侧多源验证有效 & \cite{barroso2015,harvey2018,moreiramuir2017} & \confmid \\
须做缩放归因 & 收益可能主要来自缩放而非符号 & \cite{kim2016,huang2020} & \confmid \\
恐慌状态过滤 & 崩溃可预测、动态动量翻倍 & \cite{danielmoskowitz2016} & \confmid \\
拥挤监控 & 高拥挤预示动量收益下降 & \cite{stein2009,loupolk2022,baltas2019} & \confmid \\
不做板块中性化 & 趋势收益来自共同方向风险 & 论文对偶四 & \confmid \\
不得逢低加仓 & 趋势失效含慢性出血 & 论文对偶六 & \confmid \\
\midrule
多窗口集成 & \emph{论文未检验} & — & \conflow \\
ATR 止损倍数 & \emph{论文未检验} & — & \conflow \\
免交易带系数 & \emph{论文未检验} & — & \conflow \\
恐慌阈值精确取值 & \emph{论文仅给定性刻画} & — & \conflow \\
容量上限 & \emph{论文注明无统一口径} & — & \conflow \\
\bottomrule
\end{tabular}
\end{table}

表中最后一组五行是本文刻意保留的“未支持清单”。它们不是可有可无的脚注，而是规则上线前的待办事项：任何标为\conflow~的条目，都必须先用自己的历史数据回测验证，再写入实盘规则。

\section{规则二：A 股截面反转}\label{sec:ashare}

\subsection{适用性依据}

论文表 6 的 A 股个股截面（XS）结果：

\begin{table}[htbp]
\centering\footnotesize
\caption{A 股个股截面反转的实测（论文表 6，成本 15 基点）}\label{tab:ashsrc}
\begin{tabular}{@{}>{\raggedright\arraybackslash}p{1.2cm}>{\raggedright\arraybackslash}p{1.7cm}>{\raggedright\arraybackslash}p{1.7cm}>{\raggedright\arraybackslash}p{1.7cm}>{\raggedright\arraybackslash}p{1.7cm}>{\raggedright\arraybackslash}p{1.7cm}>{\raggedright\arraybackslash}p{1.7cm}@{}}
\toprule
\textbf{$L$} & \textbf{毛利夏普} & \textbf{净夏普} & \textbf{年化收益} & \textbf{最大回撤} & \textbf{年化换手} & \textbf{月胜率} \\
\midrule
5   & $+1.08$ & $-0.35$ & $-3.26\%$ & $-39.27\%$ & $79.80$ & $46.7\%$ \\
10  & $+1.10$ & $+0.10$ & $+0.47\%$ & $-27.80\%$ & $57.74$ & $52.7\%$ \\
20  & $+1.46$ & $+0.76$ & $+6.86\%$ & $-20.53\%$ & $42.77$ & $54.7\%$ \\
60  & $+1.46$ & $+0.99$ & $+8.98\%$ & $-18.64\%$ & $28.09$ & $60.0\%$ \\
120 & $+1.24$ & $+0.86$ & $+7.17\%$ & $-22.95\%$ & $21.55$ & $57.8\%$ \\
250 & $+0.86$ & $+0.56$ & $+4.41\%$ & $-25.41\%$ & $16.37$ & $54.6\%$ \\
\bottomrule
\end{tabular}
\end{table}

四点必须读出。

\textbf{其一，信号在短窗口上很强但被成本摧毁}。$L=5$ 的毛利夏普 $1.08$，是全部实验中最强的信号之一；但 $79.80$ 倍换手带来 $12.31\%$ 的年化成本拖累，净夏普变成 $-0.35$。这是铁律二的原型案例。

\textbf{其二，$L=60$ 是全面最优}。净夏普 $0.99$、年化 $8.98\%$、最大回撤 $-18.64\%$（全窗口最浅）、月胜率 $60\%$（全窗口最高）。成本的下降（$4.69\%$）而不是信号的增强（毛利夏普与 $L=20$ 同为 $1.46$）是它胜出的原因。

\textbf{其三，优势随 $L$ 呈倒 U 形}。毛利夏普在 $L=20$ 与 $L=60$ 达到峰值 $1.46$，在两端下降（$L=5$ 的 $1.08$、$L=250$ 的 $0.86$）。这说明 A 股的反转是\emph{中期}现象，不是微观结构的买卖价差反弹——后者应当在极短窗口最强。这个判断直接影响规则设计：不应采用高频反转，而应采用月度或季度的低频反转。

\textbf{其四，做空不可实现}。论文明确披露 A 股融券受限、成本高、可借券池小，截面多空组合在多数时段不可实现，其结果应视为上限估计。这决定了本规则必须给出纯多头版本。

补充一条背景：反转策略的长期存在性有深厚的行为金融解释（过度反应后的回归\cite{debondt1985}、价值策略的外推偏差\cite{lakonishok1994}），而 A 股作为散户占比高、情绪波动大的市场，理论上更有利于这类现象的存续。但这是机制层面的推理，不构成对收益水平的承诺——论文对机制归属的置信度评级本身是中低（第~\confmid~至~\conflow~档）。

\subsection{信号与因子剥离}

\textbf{基础信号}：

\begin{equation}
s_i=\sum_{k=0}^{L-1}\ln\left(1+r_{i,k}\right) \quad\text{（过去 $L$ 个交易日累计对数收益；首选 $L=250$、备选 $60$）}.
\end{equation}

\textbf{因子剥离（必需步骤）}。论文指出行业动量解释了个股动量的大部分\cite{moskowitzgrinblatt1999}。若不剥离行业暴露，按 $s_i$ 排序得到的组合实际上在押注行业轮动，而非个股反转——这会让“反转”这一叙事失效，也会让回撤集中在少数行业。

剥离方法为横截面回归取残差（沿用论文的工程约定）：

\begin{equation}
s_i=\alpha+\sum_g\beta_g\,\mathbf 1_{\{i\in g\}}+\gamma\ln(\mathrm{MV}_i)+\varepsilon_i,
\qquad \tilde s_i=\hat\varepsilon_i ,
\end{equation}

其中 $\mathbf 1_{\{i\in g\}}$ 为行业哑变量（建议采用申万一级或中信一级），$\mathrm{MV}_i$ 为流通市值。用残差 $\tilde s_i$ 排序。

\textbf{市值中性化的理由}。小市值股票的反转幅度天然更大（波动更高、流动性更差）。若不做市值中性化，组合会系统性偏向小市值，其收益中混入规模因子的暴露，无法归因。

\subsection{组合构建}

\subsubsection*{版本 A：可执行（纯多头 + 指数对冲）}

这是唯一可实盘的版本。

\begin{itemize}[leftmargin=2.2em]
  \item \textbf{多头}：$\tilde s_i$ 最低的 20\%（即“输家”）；
  \item \textbf{加权}：等权，或按波动率倒数加权（后者可降低组合波动，但会增加小市值暴露，需配合市值中性化使用）；
  \item \textbf{对冲}：用 IC/IF/IM 空头对冲组合 beta，对冲比例 = 组合对对应指数的滚动 beta（60 日窗口，每 21 日更新）；
  \item \textbf{持仓数}：建议 $30\sim50$ 只。论文实验按 20\% 分位取，5083 只股票的池子在样本后期会给出上千只股票的组合，实盘不可行，必须限制持仓数并相应提高分位门槛。
\end{itemize}

对冲的必要性来自实证而非偏好：不对冲的纯多头反转组合的收益会与 A 股 beta 高度绑定，其最大回撤将远超表中 $-18.64\%$ 的水平（因为表中的多空组合已隐含对冲了市场方向）。

\subsubsection*{版本 B：研究（多空，明确不可实盘）}

按 $\tilde s_i$ 排序，做多最低 20\%、做空最高 20\%，两端等权。此版本用于评估信号强度与对照论文表 6，\textbf{不得用于实盘}，也不得用其实测收益做任何资金规划。

\subsection{仓位与缩放：不适用波动率目标化}

这是本规则与规则一最重要的差异，也是论文执行章的核心判断之一。

论文指出：波动率目标化在趋势策略上被反复验证有效，但把它机械照搬到反转策略上是错的——原因是信号与波动率的关系在两策略中符号相反。趋势策略在波动率上升时信号质量下降（趋势被噪声打断），降敞口是保护；而\textbf{反转策略的利润经常恰好来自高波动时段}（偏离更大，回归的绝对空间也更大）。在波动率上升时降低反转敞口，等于在信号最强时把仓位压到最小。

因此本规则明确\textbf{不做}波动率缩放。波动率信息的用途改为\emph{风险预算}与\emph{防线触发}：

\begin{itemize}[leftmargin=2.2em]
  \item 波动率显著上升（超过 80 分位）时：收紧单票止损线（由 $15\%$ 收紧至 $10\%$），检查资金冗余，但\emph{不}降低目标仓位；
  \item 波动率回落时：恢复常规止损线。
\end{itemize}

这一设计在论文中属机制推理而非直接实证（论文将“反转侧波动率目标有害”标为置信度中，缺少反转侧的系统实证）。因此本文把它标为\conflow~到\confmid~之间，并要求上线前用自有数据做对照测试。

\subsection{进出场纪律}

\textbf{调仓频率与窗口}：每 $21$ 个交易日（月频）。窗口选择与论文口径不同：论文下 A 股反转的最优区间为 $L=20\sim120$、$L=60$ 最优（月频与 $L=60$ 组合的成本占毛利约 $31.9\%$）；但在本规则的实际口径（真实申万行业中性化、纯多头、含止损与回撤控制）下重做 $L$ 扫描后，\textbf{$L=250$ 的风险特征优于 $L=60$}——两者夏普几乎相同（$0.18$ 对 $0.19$），而 $L=250$ 的最大回撤浅 $16$ 个百分点、换手低 $32\%$、成本占毛利低 $6.6$ 个百分点。因此推荐 $L=250$（首选）与 $L=60$（备选）；$L=20$ 不通过成本门槛、$L=120$ 收益为负，均排除（第~\ref{sec:retune}~节）。

\textbf{为何不用周频}。论文 $L=5$ 与 $L=10$ 的实证给出答案：$L=10$ 成本占毛利 $91.1\%$（毛利夏普 $1.10$ 对净 $0.10$），$L=5$ 直接把净夏普压成负数。周频在 A 股不可行。

\textbf{建仓与调仓细节}：
\begin{itemize}[leftmargin=2.2em]
  \item 分 $2\sim3$ 个交易日完成调仓，降低单日冲击；
  \item 涨跌停当日不成交，顺延至下一可交易日；
  \item 停牌股保留至复牌后第二个交易日再处理；
  \item 新纳入股票需满足流动性门槛（20 日均成交额 $\ge5000$ 万元）与上市满 250 日。
\end{itemize}

\subsection{失效防线}

\subsubsection*{防线一：单票偏离止损（对不收敛）}

反转策略的核心风险是“价差不收敛且继续扩大”。论文引用 Leung 与 Li(2015) 把止损退出纳入最优均值回归交易框架，正是因为它无法靠仓位缩放消除\cite{leungli2015}。

\begin{itemize}[leftmargin=2.2em]
  \item 从建仓加权成本下跌超过 $15\%$：该票止损；
  \item 跌破 60 日最低价且创出 120 日新低：该票止损（加速下跌不宜继续等待）；
  \item 高波动状态下止损线收紧至 $10\%$。
\end{itemize}

\textbf{不得逢跌加仓}。这是本规则最重要的一条纪律，也最容易违反。论文在处理 LTCM 与 2007 年量化危机时指出，反转策略的亏损本身不可怕，可怕的是亏损触发资金约束；业绩与资金流的正反馈会让套利者在错价最深时被迫减仓\cite{shleifervishny1997,khandanilo2011,brunnermeier2009}。因此加仓必须绑定在“价差已开始收敛”的证据上（例如该票已从低点反弹 5\%），而不是绑在“跌幅更大了”上。

\subsubsection*{防线二：资金冗余}

\begin{itemize}[leftmargin=2.2em]
  \item 常年预留 $30\%$ 现金，不参与建仓；
  \item 组合最大回撤达 $20\%$：目标敞口减半；
  \item 组合最大回撤达 $30\%$：全部平仓，暂停 40 个交易日。
\end{itemize}

$30\%$ 的现金冗余不是保守，而是对论文所述失效机制的针对性设计——反转策略的致命处不在亏损本身，而在被迫平仓。

\subsubsection*{防线三：行业暴露监控}

即便已做行业中性化，实际持仓的行业偏离仍会因个股停牌、涨跌停等原因累积。监控项：实际行业权重与基准行业权重的偏离度，单行业偏离超过 $5\%$ 时调整。

\subsection{成本门槛与容量}

\begin{table}[htbp]
\centering\footnotesize
\caption{A 股反转各窗口的成本门槛检查（论文口径，单边 15 基点）}\label{tab:ashcost}
\begin{tabular}{@{}>{\raggedright\arraybackslash}p{1.6cm}>{\raggedright\arraybackslash}p{2.6cm}>{\raggedright\arraybackslash}p{2.6cm}>{\raggedright\arraybackslash}p{3.6cm}@{}}
\toprule
\textbf{$L$} & \textbf{成本占毛利} & \textbf{毛利/成本} & \textbf{门槛判定（要求 $>3$）} \\
\midrule
250 & $34.5\%$ & $2.9$ & 勉强通过（临界） \\
120 & $30.7\%$ & $3.3$ & 通过 \\
60  & $31.9\%$ & $3.1$ & 通过 \\
20  & $47.6\%$ & $2.1$ & \textbf{不通过} \\
10  & $91.1\%$ & $1.1$ & \textbf{不通过} \\
5   & $132.7\%$ & $0.8$ & \textbf{不通过} \\
\bottomrule
\end{tabular}
\end{table}

门槛检查把可用区间收窄到 $L\ge60$。这一结论与论文的实操判断一致：\emph{在 A 股做截面反转，中期窗口是唯一可行区间}，短窗口不可行。

\textbf{成本上升的敏感性}。若实际成本升至 20 基点（加入冲击成本，这对大资金是现实的），成本拖累按换手比例放大，$L=60$ 的成本占毛利将从 $31.9\%$ 升至约 $42.5\%$，毛利/成本降至约 $2.4$，跌破门槛。此时须切换到 $L=120$（换手 $21.55$，成本占毛利将从 $30.7\%$ 升至约 $40.9\%$，毛利/成本约 $2.4$）——注意这仍然跌破门槛。\textbf{结论是：当单边成本超过约 18 基点时，本规则在铁律二下不再成立}。这是本规则最硬的边界。

\textbf{容量约束}。约束来自成本（论文对偶五：反转的约束是自己付的成本）。参考量级：按 $L=60$ 的 $28.09$ 倍年化换手，若组合规模为 $Q$，则年成交额约 $28Q$；在日均成交额门槛 5000 万元的股票上，需保证单票单日委托不超过其成交额的一定比例（建议 $\le5\%$）。

\subsection{参数边界表}

\begin{table}[htbp]
\centering\footnotesize
\caption{规则二参数边界}\label{tab:ashparam}
\begin{tabular}{@{}>{\raggedright\arraybackslash}p{2.6cm}>{\raggedright\arraybackslash}p{2.2cm}>{\raggedright\arraybackslash}p{3.2cm}>{\raggedright\arraybackslash}p{6.4cm}@{}}
\toprule
\textbf{参数} & \textbf{取值} & \textbf{证据来源} & \textbf{越界后果} \\
\midrule
窗口 $L$ & $250$（首选）/ $60$（备选） & 第~\ref{sec:retune}~节 & 短于 $60$ 成本占毛利超 $33\%$ 须废弃；$20$ 不通过门槛、$120$ 收益为负 \\
调仓频率 & 21 交易日 & 论文表 6 的 $L$ 区间 & 周频（$L=5/10$）毛利/成本低于 1.1，不可行 \\
分位门槛 & 20\%（实盘收窄至 $30\sim50$ 只） & 论文设定 + 工程约束 & 分位过宽则持仓过千只，不可执行 \\
因子剥离 & 行业 + log 市值 & \cite{moskowitzgrinblatt1999} + 工程约定 & 不剥离则实为押注行业轮动 \\
波动率缩放 & \textbf{不用} & 论文执行章（机制推理） & 使用则在高波动期削减最强信号 \\
单票止损线 & $15\%$（高波动 $10\%$） & \cite{leungli2015} + 工程标定 & 过紧则频繁止损；过松则失去保护 \\
现金冗余 & $30\%$ & \cite{shleifervishny1997,khandanilo2011} & 不足则可能被强制平仓 \\
回撤止损线 & $20\%$ / $30\%$ & 论文实测回撤 $\le18.64\%$ & 过紧则频繁降仓 \\
\bottomrule
\end{tabular}
\end{table}

\subsection{证据映射表}

\begin{table}[htbp]
\centering\footnotesize
\caption{规则二：决策与论文证据的映射}\label{tab:ashmap}
\begin{tabular}{@{}>{\raggedright\arraybackslash}p{3.4cm}>{\raggedright\arraybackslash}p{5.3cm}>{\raggedright\arraybackslash}p{2.6cm}>{\raggedright\arraybackslash}p{2.6cm}@{}}
\toprule
\textbf{决策} & \textbf{论文结论} & \textbf{证据来源} & \textbf{置信度} \\
\midrule
用截面反转方向 & A 股截面反转 $L=20\sim250$ 净夏普为正 & 论文表 6 & \confhigh \\
取 $L=250$（备选 $60$） & $L=60$ 净夏普 $0.99$、月胜率 $60\%$（论文口径）；$L=250$ 回撤更浅、成本更低 & 论文表 6 + 第~\ref{sec:retune}~节 & \confhigh \\
月频而非周频 & $L=5$ 毛利夏普 $1.08$ 被成本压至 $-0.35$ & 论文表 6 & \confhigh \\
必须做因子剥离 & 行业动量解释个股动量的大部分 & \cite{moskowitzgrinblatt1999} & \confmid \\
纯多头加对冲 & A 股做空不可实现（论文已披露） & 论文审计章 & \confhigh \\
不用波动率缩放 & 反转利润常来自高波动时段 & 论文执行章 & \confmid \\
单票偏离止损 & 止损须纳入最优反转交易框架 & \cite{leungli2015} & \confmid \\
资金冗余 & 资金约束是反转策略的致命处 & \cite{shleifervishny1997,khandanilo2011} & \confmid \\
不得逢跌加仓 & 亏损触发资金约束即被迫平仓 & \cite{brunnermeier2009} & \confmid \\
\midrule
对冲 beta 的窗口长度 & \emph{论文未检验} & — & \conflow \\
分仓完成调仓的交易日数 & \emph{论文未检验} & — & \conflow \\
行业分类标准选择 & \emph{论文未检验} & — & \conflow \\
\bottomrule
\end{tabular}
\end{table}

\section{规则三：期货截面统计套利}\label{sec:statarb}

\subsection{适用性依据}

论文表 5 的期货截面口径（XS）结果：

\begin{table}[htbp]
\centering\footnotesize
\caption{期货截面反转的实测（论文表 5，成本 2 基点）}\label{tab:stasrc}
\begin{tabular}{@{}>{\raggedright\arraybackslash}p{1.2cm}>{\raggedright\arraybackslash}p{1.7cm}>{\raggedright\arraybackslash}p{1.7cm}>{\raggedright\arraybackslash}p{1.7cm}>{\raggedright\arraybackslash}p{1.7cm}>{\raggedright\arraybackslash}p{1.7cm}>{\raggedright\arraybackslash}p{1.5cm}@{}}
\toprule
\textbf{$L$} & \textbf{毛利夏普} & \textbf{净夏普} & \textbf{年化收益} & \textbf{最大回撤} & \textbf{年化换手} & \textbf{NW $t$} \\
\midrule
5   & $-0.05$ & $-0.20$ & $-2.59\%$ & $-33.53\%$ & $79.58$ & $-0.76$ \\
10  & $-0.20$ & $-0.31$ & $-3.81\%$ & $-46.20\%$ & $57.81$ & $-1.19$ \\
20  & $+0.07$ & $-0.01$ & $-0.71\%$ & $-34.84\%$ & $41.76$ & $-0.04$ \\
60  & $+0.77$ & $+0.72$ & $+7.57\%$ & $-18.74\%$ & $26.63$ & $+2.59$ \\
120 & $+0.60$ & $+0.56$ & $+5.44\%$ & $-17.40\%$ & $20.78$ & $+1.98$ \\
250 & $+0.28$ & $+0.25$ & $+2.09\%$ & $-26.64\%$ & $16.08$ & $+0.82$ \\
\bottomrule
\end{tabular}
\end{table}

论文表 5 中的这组结果有一个值得单独强调的地位：$L=60$ 的 Newey--West $t$ 为 $2.59$，是论文全部 48 个组合中少数通过常规显著性门槛的结果。同时它的最大回撤仅 $-18.74\%$、成本占毛利约 $6.3\%$（在全部门槛检查中余量最充足）。

三点必须读出。

\textbf{其一，短窗口在这条口径上是负收益}。$L=5$ 与 $L=10$ 的毛利夏普已经是负的（$-0.05$、$-0.20$），$L=20$ 接近零（$+0.07$）。这与论文的 U 形结构判断一致：期货截面上的短线趋势（$L\le10$）与中线反转（$L\ge60$）之间存在一个过渡带，$L=20$ 恰在过渡带中央，两边都不占优。

\textbf{其二，$L=60$ 与 $L=120$ 构成可行区间}。二者的净夏普 $0.72$ 与 $0.56$、回撤 $-18.74\%$ 与 $-17.40\%$、$t$ 值 $2.59$ 与 $1.98$。$L=250$ 的收益与显著性都明显衰减。

\textbf{其三，本规则的实现方式与论文不同，须明确标注}。论文的截面口径是“按过去 $L$ 日累计收益排序，多后 20\%、空前 20\%”，这是一个\emph{纯截面排序}策略，不使用协整、不估计 OU、不做对对冲。而本文设计的规则三走向\emph{协整配对}路线，理由有两条：其一，纯截面排序在多空两端会引入大量方向性风险（组合对商品指数有净暴露），而配对天然对冲共同因子；其二，论文的对偶七指出反转策略在结构上应带“资格检验”，协整检验正是这道门槛，可以让规则在最弱的组合（$L=20$ 附近的过渡带）上直接弃权。

必须诚实说明：\textbf{配对版本的历史绩效不能直接引用论文表 5 的数字}——两者是不同的策略构造。论文表 5 支持的是“期货截面上中线反转方向有效”这一判断，本文据此选择 $L=60/120$ 作为信号窗口，但配对化的收益水平需自行回测，本文不给出预期数字。

\subsection{配对与篮子形成}

配对交易的方法论基础由 Vidyamurthy(2004) 的专著系统化，其覆盖的距离法、协整法与随机价差法三种配对选择思路至今仍是该领域的技术主干\cite{vidyamurthy2004}；Krauss(2017) 的综述梳理了该族策略的方法演进与利润衰减脉络\cite{krauss2017}。本规则的构造（产业链内配对 + 协整筛选 + OU 阈值）属于其中的协整分支。

\textbf{配对范围：产业链内}。跨产业链配对在商品期货上往往不成立（例如黑色与贵金属之间缺乏经济意义上的长期均衡关系）。建议的产业链划分：

\begin{table}[htbp]
\centering\footnotesize
\caption{建议的配对产业链划分}\label{tab:chains}
\begin{tabular}{@{}>{\raggedright\arraybackslash}p{3.2cm}>{\raggedright\arraybackslash}p{10.8cm}@{}}
\toprule
\textbf{产业链} & \textbf{候选品种} \\
\midrule
黑色系 & RB（螺纹）、HC（热卷）、I（铁矿）、J（焦炭）、JM（焦煤）、SF（硅铁）、SM（锰硅） \\
能源化工 & MA（甲醇）、PP（聚丙烯）、L（塑料）、TA（PTA）、EG（乙二醇）、V（PVC）、BU（沥青）、SC（原油）、FU（燃油） \\
有色金属 & CU（铜）、AL（铝）、ZN（锌）、NI（镍）、SN（锡）、PB（铅） \\
农产品 & M（豆粕）、Y（豆油）、P（棕榈）、C（玉米）、CS（淀粉）、OI（菜油）、RM（菜粕）、SR（白糖）、CF（棉花）、AP（苹果） \\
油脂油料 & A（豆一）、B（豆二）、M、Y、P、OI \\
\bottomrule
\end{tabular}
\end{table}

\textbf{形成期与交易期}。形成期 60 交易日，交易期 21 交易日，滚动推进。这一“形成期与交易期严格分离”的设计来自 Gatev 等(2006) 的 6+6 个月滚动范式\cite{gatev2006}，其目的是消除 look-ahead——若配对是在包含交易期的数据上选出的，价差平稳性会被机械地保证。

\textbf{候选对生成}。同一产业链内两两组合，形成期相关性 $\ge0.6$ 者进入候选池。

\subsection{协整资格检验（入场门槛）}

这道门槛是本规则区别于规则一、二的核心结构，其依据是论文对偶七：反转策略在结构上多一道“资格检验”，是“条件性”策略。

\textbf{两步法}（Engle--Granger\cite{englegranger1987}）：

\begin{enumerate}[leftmargin=2.2em]
  \item 对形成期数据做回归 $\ln P^A_t=\alpha+\beta\ln P^B_t+e_t$；
  \item 对残差 $e_t$ 做 ADF 检验\cite{dickeyfuller1979}，统计量需小于 5\% 临界值。
\end{enumerate}

不通过的对直接弃权，不进入交易。这道门槛的价值不在于提高单对收益，而在于\emph{减少无效样本}——论文指出反转策略的“资格检验”会消耗样本，使有效样本基础比趋势策略更薄，这是它的结构性代价。

\subsection{OU 参数估计与半衰期筛选}

对通过协整检验的价差 $e_t$ 拟合 AR(1)：

\begin{equation}
e_{t+1}=c+\phi e_t+u_t \;\Longrightarrow\;
\theta=-\frac{\ln\phi}{\Delta},\qquad
H=\frac{\ln 2}{\theta} .
\end{equation}

\textbf{半衰期筛选区间}：$H\in[5,\,60]$ 交易日。上界的含义是资金占用：$H$ 过大意味着价差需要数月才收敛一半，占用保证金的效率太低；下界的含义是成本：$H$ 过小意味着价差变化极快，交易频率上升、成本侵蚀收益。

\textbf{估计偏差的提醒}。论文附录指出 $\hat\phi$ 在小样本下有向下偏误，短样本估出的 $\theta$ 往往偏大、半衰期偏短\cite{kendall1953}。形成期 60 交易日对 AR(1) 估计而言偏短，因此建议对 $\hat\phi$ 做偏误修正，或把形成期延长至 120 交易日。这一点在本文中属\textbf{工程判断}。

\subsection{对冲比率}

\textbf{静态版本的缺陷}。若对冲比率 $\beta$ 在整个交易期保持不变（用形成期的 OLS 系数），则当两品种的相对关系发生漂移时，价差序列会失真——它会把参数漂移误读为价差偏离，产生虚假信号。

\textbf{在线版本（推荐）}。用 Kalman 滤波滚动估计 $\beta_t$，其状态空间形式为

\begin{equation}
\begin{aligned}
\ln P^A_t &= \alpha_t+\beta_t\ln P^B_t+\varepsilon_t, &&\varepsilon_t\sim\mathcal N(0,\sigma_\varepsilon^2),\\
\alpha_{t+1}&=\alpha_t+\eta^{\alpha}_t, \quad
\beta_{t+1}=\beta_t+\eta^{\beta}_t, &&\eta\sim\mathcal N(0,\Sigma_\eta),
\end{aligned}
\end{equation}

用时刻 $t$ 及之前的信息更新 $(\alpha_t,\beta_t)$，据此计算价差 $e_t=\ln P^A_t-\alpha_t-\beta_t\ln P^B_t$。这一做法来自 Elliott 等(2005) 的 OU 状态空间配对框架\cite{elliott2005}。

\subsection{交易带（本规则的核心）}

\textbf{标准化价差}：

\begin{equation}
z_t=\frac{e_t-\mu_t^{(60)}}{\sigma_t^{(60)}} ,
\end{equation}

其中 $\mu_t^{(60)}$ 与 $\sigma_t^{(60)}$ 为价差的滚动 60 日均值与标准差。

\textbf{带宽的理论依据}。论文附录给出：最优无交易带的宽度满足

\begin{equation}
b\;\propto\;\left(\frac{\text{单位交易成本}}{\text{风险厌恶}\times\theta}\right)^{1/3},
\end{equation}

即带宽随成本以三次方根增长\cite{martin2011}。这条关系是本文“反转策略不能频繁穿越阈值”这一工程要求的理论出处。

\textbf{离散化的事件规则}（把连续带宽翻译成可执行的进出场）：

\begin{table}[htbp]
\centering\footnotesize
\caption{期货统计套利的进出场规则}\label{tab:band}
\begin{tabular}{@{}>{\raggedright\arraybackslash}p{2.4cm}>{\raggedright\arraybackslash}p{4.0cm}>{\raggedright\arraybackslash}p{7.6cm}@{}}
\toprule
\textbf{条件} & \textbf{动作} & \textbf{依据} \\
\midrule
$|z_t|>2.0$，无持仓 & 开仓：$z>2$ 空 A 多 B；$z<-2$ 反之 & s-score 进出场范式\cite{avellaneda2010} \\
$|z_t|<0.5$，有持仓 & 平仓 & 同上 \\
$|z_t|>3.5$，有持仓 & 强制止损 & 不收敛防线\cite{leungli2015} \\
持仓 $>3H$ 且未收敛 & 平仓 & 时间止损（资金效率） \\
\bottomrule
\end{tabular}
\end{table}

\textbf{参数边界}。$2.0$ 与 $0.5$ 取自 Avellaneda--Lee 的 s-score 范式（开仓阈值 $2$、平仓阈值 $0.5\sim0.75$）\cite{avellaneda2010}；$3.5$ 的止损阈值为工程标定。按三次方根关系，若实际单边成本从 2 基点升至 6 基点（3 倍），带宽应约放宽至 $3^{1/3}\approx1.44$ 倍，即开仓阈值由 $2.0$ 提高到约 $2.9$。这是本文给出的唯一一条\emph{有理论形式}的参数调整规则，值得在实盘中优先采用。

\subsection{组合与分散}

\textbf{持仓数}：建议同时持有 $N\ge10$ 对。论文指出配对交易的价差交易天然对冲了共同因子暴露，但代价是特质风险无法被市场因子抵消，必须靠配对数量分散（论文执行章）。

\textbf{权重}：按价差波动率倒数归一，使每对的风险贡献大致相等。

\textbf{不做波动率目标化}。与规则二同理：反转类策略不用波动率缩放仓位。论文执行章明确指出对反转策略应把波动率信息用于风险预算（收紧止损、检查冗余），而非信号缩放。

\subsection{失效防线}

\subsubsection*{防线一：不收敛止损}

\begin{itemize}[leftmargin=2.2em]
  \item $|z_t|>3.5$：该对强制平仓；
  \item 持仓超过 $3H$ 仍未回到 $|z|<0.5$：该对平仓；
  \item 单对累计亏损超过该对风险预算的 2 倍：平仓。
\end{itemize}

\subsubsection*{防线二：组合止损与资金冗余}

\begin{itemize}[leftmargin=2.2em]
  \item 组合最大回撤达 $20\%$：敞口减半；
  \item 组合最大回撤达 $30\%$：全部平仓，暂停 40 个交易日；
  \item 预留保证金冗余不低于名义敞口的 $40\%$（期货杠杆下这条尤其重要）。
\end{itemize}

\subsubsection*{防线三：产业链集中度}

同一产业链内的持仓对数不超过总对数的一半。理由是产业链内的品种高度相关，若某一产业链的均衡关系整体失效（如政策冲击），集中持仓会同时受损。

\subsubsection*{共同失效的警告}

论文对偶六的边界条目指出：\textbf{两策略的失效在极端去杠杆事件中会同时发生}——2007 年 8 月中，量化反转基金与动量基金都被冲击，因为冲击来自共同的资金端而非各自的信号端\cite{khandanilo2011,brunnermeier2009}。这意味着若同时运行规则一与规则三，“分散化”在资金冲击下会失效。组合层面必须为此预留额外冗余，而不能假设两个策略互为对冲。

\subsection{成本门槛与容量}

\textbf{成本核算}。单边 2 基点，但必须加上一条修正：\emph{反转类交易的冲击成本高于平均值}，因为交易密集发生在流动性较差的时点。论文审计章把这一点列为已知的未修正偏差。实盘应切换到参与率约束的平方根冲击模型。

\textbf{门槛检查}。若以论文表 5 的 $L=60$ 数据作参照（成本占毛利 $6.3\%$，毛利/成本约 $15.9$），本规则的成本余量在全部门槛检查中最充足。但如前所述，配对版本的实际成本占比需自行回测确定。

\textbf{容量约束}。约束来自两处：其一是成本；其二是产业链内的品种深度——同一产业链内可用于配对的流动性规模有限，且配对交易在两个方向同时占用保证金。

\subsection{参数边界表}

\begin{table}[htbp]
\centering\footnotesize
\caption{规则三参数边界}\label{tab:staparam}
\begin{tabular}{@{}>{\raggedright\arraybackslash}p{2.6cm}>{\raggedright\arraybackslash}p{2.2cm}>{\raggedright\arraybackslash}p{3.2cm}>{\raggedright\arraybackslash}p{6.4cm}@{}}
\toprule
\textbf{参数} & \textbf{取值} & \textbf{证据来源} & \textbf{越界后果} \\
\midrule
信号窗口 & 60（备选 120） & 论文表 5 & 短于 20 毛利夏普已为负；$L=20$ 处于过渡带 \\
形成期 / 交易期 & 60 / 21 交易日 & \cite{gatev2006} 滚动范式 & 重叠则价差平稳性被机械保证（look-ahead） \\
协整显著性 & ADF $<5\%$ 临界值 & \cite{englegranger1987,dickeyfuller1979} & 放宽则无效配对涌入 \\
半衰期区间 & $[5,60]$ 交易日 & 论文附录 A.1 & 下限过低成本激增；上限过高资金效率低 \\
开仓 / 平仓阈值 & $2.0$ / $0.5$ & \cite{avellaneda2010} & 阈值过低频繁交易；过高错过机会 \\
止损阈值 & $3.5$ & 工程标定 & 过宽则单对亏损失控 \\
成本上调的带宽调整 & $b\propto c^{1/3}$ & \cite{martin2011} & 成本涨 3 倍须把阈值提至约 $2.9$ \\
持仓对数 & $\ge10$ & 论文执行章 & 过少则特质风险无法分散 \\
产业链集中度 & $\le50\%$ & 工程判断 & 过度集中则单一产业链冲击致命 \\
保证金冗余 & $\ge40\%$ & \cite{brunnermeier2009} & 不足则可能被强制平仓 \\
\bottomrule
\end{tabular}
\end{table}

\subsection{证据映射表}

\begin{table}[htbp]
\centering\footnotesize
\caption{规则三：决策与论文证据的映射}\label{tab:stamap}
\begin{tabular}{@{}>{\raggedright\arraybackslash}p{3.4cm}>{\raggedright\arraybackslash}p{5.3cm}>{\raggedright\arraybackslash}p{2.6cm}>{\raggedright\arraybackslash}p{2.6cm}@{}}
\toprule
\textbf{决策} & \textbf{论文结论} & \textbf{证据来源} & \textbf{置信度} \\
\midrule
期货截面反转方向有效 & $L=60/120$ 净夏普 $0.72/0.56$、$t=2.59/1.98$ & 论文表 5 & \confmid \\
取 $L=60$ 为主窗口 & 全部门槛检查中成本余量最充足 & 论文表 5 & \confmid \\
弃用 $L\le20$ & 毛利夏普已为负或接近零 & 论文表 5 & \confhigh \\
必须做协整资格检验 & 反转类型是条件性策略 & 论文对偶七 & \confmid \\
半衰期筛选 & OU 离散估计与半衰期定义 & 论文附录 A.1 & \confmid \\
带宽随成本立方根变化 & 最优无交易带宽度关系 & \cite{martin2011} & \confmid \\
在线估计对冲比率 & OU 状态空间配对框架 & \cite{elliott2005} & \confmid \\
s-score 进出场范式 & 因子剥离后的统计套利实现 & \cite{avellaneda2010} & \confmid \\
不用波动率缩放 & 同规则二 & 论文执行章 & \confmid \\
止损与资金冗余 & 不收敛加资金断裂是致命处 & \cite{leungli2015,brunnermeier2009} & \confmid \\
\midrule
配对化的收益水平 & \emph{论文未检验配对版本，不可引用表 5 数字} & — & \conflow \\
开仓阈值 $2.0$ 的具体取值 & \emph{论文未给本文样本下的标定} & — & \conflow \\
对冲比 Kalman 参数 & \emph{论文未检验} & — & \conflow \\
产业链划分方案 & \emph{论文未检验} & — & \conflow \\
\bottomrule
\end{tabular}
\end{table}

表中第一行给 \confmid~而非 \confhigh，反映一个必须正视的事实：论文表 5 支持的是“纯截面排序”形式的反转，而本规则把它改造成了“协整配对”形式。这一步改造有理论依据（对偶七、Avellaneda--Lee 范式），但缺少本文样本下的直接实证。这是本规则与规则二在证据强度上的实质差别。

\section{规则对比与选用}\label{sec:compare}

\subsection{三套规则的横向对照}

\begin{table}[htbp]
\centering\footnotesize
\caption{三套规则对照}\label{tab:compare}
\begin{tabular}{@{}>{\raggedright\arraybackslash}p{2.9cm}>{\raggedright\arraybackslash}p{3.7cm}>{\raggedright\arraybackslash}p{3.7cm}>{\raggedright\arraybackslash}p{3.5cm}@{}}
\toprule
\textbf{维度} & \textbf{规则一 期货趋势} & \textbf{规则二 A 股反转} & \textbf{规则三 期货统计套利} \\
\midrule
信号口径 & 时序（单品种对自身历史） & 截面（全市场排序） & 截面（产业链内配对） \\
方向 & 趋势（$+s$） & 反转（$-s$） & 反转（价差偏离） \\
主窗口 & $L=250$ / 60 / 10 & $L=60$（备选 120） & $L=60$（备选 120） \\
调仓频率 & 5 交易日 & 21 交易日 & 事件驱动（阈值触发） \\
年化换手 & $8.8\sim35$ & $28.1$ & $\sim27$（论文参照值） \\
成本占毛利 & $10.2\%\sim29.4\%$ & $31.9\%$ & $\sim6.3\%$（论文参照值） \\
波动率目标化 & \textbf{用} & 不用 & 不用 \\
入场门槛 & 信号强度 $>1\sigma$ & 无（排序） & \textbf{协整检验 + 半衰期筛选} \\
因子中性化 & 不做 & \textbf{行业 + 市值} & 配对天然对冲 \\
主要失效模式 & 慢性出血 + 崩溃跳崖 & 不收敛 + 资金断裂 & 不收敛 + 保证金挤兑 \\
主要防线 & 恐慌状态 + 拥挤监控 & 单票止损 + 现金冗余 & 止损 + 保证金冗余 \\
容量约束来源 & 拥挤（信号同质化） & 成本 & 成本 + 产业链深度 \\
论文支持的置信度 & 中（表 5 时序口径直接支持） & 高（表 6 直接支持） & 中（方向支持，实现改造未经检验） \\
\bottomrule
\end{tabular}
\end{table}

\emph{注}：表中“主窗口”一行的取值沿用论文口径；第~\ref{sec:retune}~节的重标定把规则一改为主推 $L=10$（备选 $20$）、规则二改为主推 $L=250$（备选 $60$），其余维度不变。“成本占毛利”一行同样来自论文口径，工程增强后的实测值见第~\ref{sec:verify}~节。

\subsection{三套规则的共同点与差异}

\textbf{唯一的形式差异是符号，其余差异全部在执行层}。这是本文（以及其依据的论文）最核心的判断。表中“信号口径”一栏看起来三种不同，但拆开看：规则一用 $+\mathrm{sign}(s)$，规则二用 $-\mathrm{sign}(s)$，规则三用 $-\mathrm{sign}(e)$，而 $e$ 本身就是两资产价格的对数差，等价于价差的累计收益。三者都是“过去价格变化的单调函数”，符号分别是 $+$、$-$、$-$。

真正的差异在四行：\emph{波动率目标化的取舍}、\emph{入场门槛的有无}、\emph{中性化的必要性}、\emph{防线的内容}。这四行恰好对应论文七组对偶中最具工程后果的四组。

\textbf{为什么规则一用波动率目标而另两套不用}。这不是风格偏好，而是信号与波动率的联合分布不同。趋势信号在波动率上升时质量下降，反转信号在高波动时反而更强（偏离更大、回归空间更大）。若对三套规则统一施加波动率目标，会系统性地削弱规则二与规则三——论文执行章称这是“在信号最强时把仓位压到最小”。

\textbf{为什么只有规则三有协整门槛}。论文对偶七指出，反转策略在结构上应带资格检验（“先验证后交易”），而趋势策略没有这道门槛（“边交易边估计”）。规则二虽然也是反转，但它走的是全市场排序，排序本身隐含了流动性筛选，没有第二道结构门槛。规则三做配对，则必须显式检验协整关系——否则“价差会回归”这一前提不成立。

\subsection{选用决策树}

\begin{enumerate}[leftmargin=2.2em]
  \item \textbf{能否做空？}
  \begin{itemize}[leftmargin=1.8em]
    \item 不能（A 股账户为主，无期货或融券）：只能选\textbf{规则二}，且必须走版本 A（纯多头 + 指数对冲）。
    \item 能（有期货账户）：继续。
  \end{itemize}
  \item \textbf{期望持仓周期？}
  \begin{itemize}[leftmargin=1.8em]
    \item 少于 5 个交易日：三套规则都不适用。论文实证显示短窗口在全部口径上都被成本摧毁（期货趋势 $L=5$ 成本占毛利 $50.1\%$；A 股反转 $L=5$ 毛利夏普 $1.08$ 被压成 $-0.35$）。
    \item 5 至 60 个交易日：继续。
  \end{itemize}
  \item \textbf{首要目标？}
  \begin{itemize}[leftmargin=1.8em]
    \item 低换手、大资金：选\textbf{规则一}的 $L=250$ 档（换手 $8.83$，成本占毛利 $10.2\%$）。
    \item 追求最高历史夏普：选\textbf{规则二}的 $L=60$ 档（净夏普 $0.99$），前提是接受 A 股做空实现与除权修复的残余偏差。
    \item 已有期货账户、希望降低商品指数方向暴露：选\textbf{规则三}，但须自行回测配对版本的收益水平。
  \end{itemize}
  \item \textbf{是否需要同时运行多套？}
  \begin{itemize}[leftmargin=1.8em]
    \item 若要同时运行规则一与规则三，必须注意论文对偶六的警告：两策略的失效在极端去杠杆事件中会同时发生（2007 年 8 月），因为冲击来自共同的资金端而非各自的信号端。此时“分散化”会失效，组合层面须预留额外冗余。
  \end{itemize}
\end{enumerate}

\emph{更正（见第~\ref{sec:retune}~节）}：第 3 步的窗口取值沿用论文口径，工程增强后的重标定结果与之不同——规则一首选 $L=10$、规则二首选 $L=250$。决策树的分支逻辑不变，仅窗口的具体取值以第~\ref{sec:retune}~节为准。

\subsection{一个必须避免的组合方式}

有一种做法看起来合理但实际有害：\textbf{把趋势与反转等权叠加，期望获得分散化}。论文执行章对此的分析是，两策略在某些状态下相关性趋近于一（共同的资金冲击），在另一些状态下相关性为负（信号符号相反）；简单叠加的实际效果是“在信号相反时互相抵消，在资金冲击时同时亏损”。

正确的融合方式有三条（均为工程建议，论文未检验）：
\begin{itemize}[leftmargin=2.2em]
  \item \textbf{按尺度分层}：短尺度配反转、中尺度配趋势，让不同频段各自发挥作用；
  \item \textbf{按状态加权}：用 regime 概率动态调整两策略的权重，而非固定比例；
  \item \textbf{按风险预算而非资金分配}：先评估两策略在目标状态下的波动与尾部，再分配风险预算。
\end{itemize}

\section{证据映射总表与上线检查单}\label{sec:mapping}

\subsection{三套规则的证据汇总}

下表把三套规则的关键决策集中列出，按“论文表 5/表 6 直接支持”、“文献支持（论文转引）”、“论文论证”、“实测修正”、“工程判断”五档标注，供上线评审时逐条核对。其中“实测修正”一档来自第~\ref{sec:verify}~节与第~\ref{sec:retune}~节对论文口径的推翻与重标定。

\begin{longtable}{@{}>{\raggedright\arraybackslash}p{2.4cm}>{\raggedright\arraybackslash}p{4.6cm}>{\raggedright\arraybackslash}p{4.4cm}>{\raggedright\arraybackslash}p{2.6cm}@{}}
\caption{三套规则的证据分级总表}\label{tab:allmap}\\
\toprule
\textbf{规则} & \textbf{决策} & \textbf{依据} & \textbf{级别} \\
\midrule
\endfirsthead
\toprule
\textbf{规则} & \textbf{决策} & \textbf{依据} & \textbf{级别} \\
\midrule
\endhead
\bottomrule
\endlastfoot
规则一 & 期货时序用趋势方向 & 论文表 5：全窗口净夏普为正 & 实证直接支持 \\
规则一 & 弃用 $L=5$ & 论文表 5：成本占毛利 $50.1\%$ & 实证直接支持 \\
规则一 & 弃用 $L=250$ & 工程增强后成本占毛利 $108.7\%$（第~\ref{sec:verify}~节） & 实测修正 \\
规则一 & 窗口取 $L=10$（首选）、$20$（备选） & 第~\ref{sec:retune}~节：参数曲面平滑、成本占毛利 $\le18\%$ & 实测修正 \\
规则一 & 波动率目标化 & \cite{barroso2015,harvey2018,moreiramuir2017} & 文献支持 \\
规则一 & 缩放归因检验 & \cite{kim2016,huang2020} & 文献支持 \\
规则一 & 恐慌状态过滤 & \cite{danielmoskowitz2016} & 文献支持 \\
规则一 & 拥挤监控 & \cite{stein2009,loupolk2022,baltas2019} & 文献支持 \\
规则一 & 不做板块中性化 & 论文对偶四（收益来自共同方向风险） & 论文论证 \\
规则一 & 免交易带系数 $0.1$ & — & 工程判断 \\
规则一 & ATR 止损倍数 & — & 工程判断 \\
规则一 & 恐慌阈值 $-5\%$/80 分位 & DM16 定性 + 自标定 & 工程判断 \\
\midrule
规则二 & A 股截面用反转方向 & 论文表 6：$L=20\sim250$ 净夏普为正 & 实证直接支持 \\
规则二 & 取 $L=250$（备选 $60$） & 第~\ref{sec:retune}~节：回撤浅 $16$ 个百分点、成本占毛利低 $6.6$ 个百分点 & 实测修正 \\
规则二 & 月频而非周频 & 论文表 6：$L=5$ 毛利夏普 $1.08$ 被压至 $-0.35$ & 实证直接支持 \\
规则二 & 纯多头 + 指数对冲 & 论文审计章：A 股做空不可实现 & 实证直接支持 \\
规则二 & 行业 + 市值中性化 & \cite{moskowitzgrinblatt1999} & 文献支持 \\
规则二 & 不用波动率缩放 & 论文执行章（机制推理） & 论文论证 \\
规则二 & 单票偏离止损 & \cite{leungli2015} & 文献支持 \\
规则二 & 现金冗余 $30\%$ & \cite{shleifervishny1997,khandanilo2011} & 文献支持 \\
规则二 & 不得逢跌加仓 & \cite{brunnermeier2009} & 文献支持 \\
规则二 & 对冲 beta 窗口长度 & — & 工程判断 \\
规则二 & 行业分类标准 & — & 工程判断 \\
\midrule
规则三 & 期货截面中线用反转 & 论文表 5：$L=60/120$ 净夏普 $0.72/0.56$、$t=2.59/1.98$ & 实证直接支持 \\
规则三 & 弃用 $L\le20$ & 论文表 5：毛利夏普为负或近零 & 实证直接支持 \\
规则三 & 配对的实现方式 & 论文对偶七 + \cite{avellaneda2010} & 论文论证（配对版收益未检验） \\
规则三 & 协整资格检验 & \cite{englegranger1987,dickeyfuller1979} & 文献支持 \\
规则三 & 半衰期筛选 $[5,60]$ & 论文附录 A.1 & 论文论证 \\
规则三 & 带宽 $\propto$ 成本立方根 & \cite{martin2011} & 文献支持 \\
规则三 & Kalman 在线对冲比 & \cite{elliott2005} & 文献支持 \\
规则三 & s-score 进出场阈值 & \cite{avellaneda2010} & 文献支持 \\
规则三 & 止损阈值 $3.5$ & — & 工程判断 \\
规则三 & 产业链划分方案 & — & 工程判断 \\
\end{longtable}

\subsection{仍未补齐的检验}

论文的实证有明确的未竟事项。本文已补做两项——多重检验校正（第~\ref{sec:mtflow}~节）与规则三的配对筛选及回测（第~\ref{sec:pairalgo}~节）；其余四项仍未完成。它们不是可选项，而是规则成立的前提。

\begin{enumerate}[leftmargin=2.2em]
  \item \textbf{涨跌停与停牌的可交易性建模}。论文未做。需要在信号与收益计算中显式排除涨跌停板与停牌日的成交，否则 2015 年等异常波动时段的可实现收益被系统性高估。本文的规则文本已写入顺延规则，但回测引擎尚未接入（第~\ref{sec:mtflow}~节第 7 步因此判为未通过）。
  \item \textbf{参与率约束的冲击成本}。论文用固定基点。反转类规则（规则二、三）的换手高且交易密集于流动性差的时点，必须改用 \mbox{impact}$=K\sqrt{\text{参与率}}$ 模型并标定 $K$。依据：\cite{almgrenchriss2000,perold1988}。
  \item \textbf{A 股复权因子的补全}。论文用“跌幅超涨跌停即判除权”的近似规则，漏掉小额分红型除权。需从数据商处取得复权因子序列替换。
  \item \textbf{规则一的缩放与信号分离归因}。分别测试“只保留符号”与“只做缩放”的对照组，确认收益来源。依据：\cite{kim2016,huang2020}。
\end{enumerate}

\subsection{持续监控清单}

规则上线后需持续监控的四类指标，以及触发时的动作：

\begin{table}[htbp]
\centering\footnotesize
\caption{持续监控清单}\label{tab:monitor}
\begin{tabular}{@{}>{\raggedright\arraybackslash}p{3.4cm}>{\raggedright\arraybackslash}p{4.6cm}>{\raggedright\arraybackslash}p{3.6cm}>{\raggedright\arraybackslash}p{2.4cm}@{}}
\toprule
\textbf{监控项} & \textbf{指标定义} & \textbf{触发阈值} & \textbf{动作} \\
\midrule
成本占比 & 实测成本占毛利比 & 规则一 $>35\%$；规则二 $>40\%$ & 切换更长窗口或停机 \\
拥挤度 & 策略内两两平均相关（滚动 3 年分位） & $>90$ 分位 & 规则一敞口 $\times0.7$ \\
恐慌状态 & 市场 60 日收益与波动率分位 & 收益 $<-5\%$ 且波动 $>$80 分位 & 规则一敞口 $\times0.5$ \\
实测偏度 & 日收益偏度 & 趋势规则转负且显著 & 复核是否已失去凸性 \\
最大回撤 & 组合净值回撤 & 规则一 $15\%/25\%$；规则二三 $20\%/30\%$ & 腰斩 / 全平 \\
资金冗余 & 现金或保证金缓冲比例 & 低于设计值 & 降低敞口而非补足杠杆 \\
行业偏离 & 实际行业权重与基准偏离 & 单行业 $>5\%$ & 调整持仓 \\
产业链集中度 & 单产业链持仓占比 & $>50\%$ & 调整配对分布 \\
\bottomrule
\end{tabular}
\end{table}

\subsection{最后一条提醒}

本文给出的三套规则都有一个共同的性质：\textbf{它们的历史绩效来自一个未做可交易性过滤、成本模型简化的对照实验}。论文已经把这些局限逐条披露（审计章），本文在规则设计层面对其中三项做了补救——可交易性处理、成本门槛、失效防线，并在第~\ref{sec:mtflow}~节补做了多重检验校正。但回到可实现性，第~\ref{sec:mapping}~节列出的四项仍未完成：涨跌停与停牌过滤尚未接入回测引擎、冲击成本模型未标定、A 股复权因子未补全、规则一的缩放归因未分离。

在补齐之前，这三套规则应当被当作\emph{设计模板}而不是\emph{可上线策略}。把模板误当作策略，正是论文反复警告的那类错误——用统计上尚不足以称为显著的结果，去支撑实际的资金决策。

\section{伪代码汇总}\label{sec:pseudocode}

本节把回测框架与三套规则的执行逻辑写成伪代码。写法遵循本文的既有约定：不引入语言特定的库依赖，不包含导入与函数定义，只描述变量、条件与动作的流转，便于直接翻译为任意实现语言。框架与策略分两层：先给三套规则共用的引擎骨架，再给各自策略层的执行逻辑。

\subsection{回测框架：引擎与策略的分层}

三套规则跑在同一个引擎上，差异全部封装在“取目标权重”这一步。引擎负责口径一致，策略负责信号正确——把两者分开，第~\ref{sec:mtflow}~节的流程才有统一的挂载点。

\begin{quote}\ttfamily\small
初始化：\\
\quad 加载数据层（第~\ref{sec:shared}~节口径：复权、换月、可交易标记）\\
\quad 加载成本参数、评估指标口径、防线阈值\\[2pt]
对每个调仓周期（频率由策略声明：规则一每 5 个交易日，规则二、三每 21 个交易日）：\\
\quad \textbf{取目标权重} $\gets$ 调用当前策略的信号逻辑（见下三小节）\\
\quad 可交易性过滤：\\
\qquad 对每个标的，若当日涨跌停或停牌：本次不下单，顺延至下一可交易日\\
\qquad 若冲击模型算出参与率 $>10\%$：标记该笔，拆单或放弃\\[2pt]
\quad \textbf{换手与成本}：\\
\qquad 换手 $\gets\sum_i\bigl|w_{i}^{\text{target}}-w_{i}^{\text{current}}\bigr|$\\
\qquad 成本 $\gets$ 仅对换入换出部分按固定基点收费 $+$ 冲击项\\
\qquad 持有未动的仓位零成本\\[2pt]
\quad \textbf{执行}：按目标权重下单；未成交部分保留旧权重，不重复收费\\[2pt]
每个交易日：\\
\quad 持仓漂移：$w_i \gets w_i\times(1+r_{i,t})/\bigl(1+\sum_j w_j r_{j,t}\bigr)$\\
\quad 调用策略的防线检查（止损、回撤降敞口、状态过滤）\\[2pt]
收尾（每次运行）：\\
\quad 输出评估指标（含毛利与成本分列）\\
\quad 记录本次运行的参数、口径版本与结果指纹，写入试验账本\\
\quad 触发上线前流程的回归测试（第~\ref{sec:mtflow}~节）
\end{quote}

两条分层纪律值得单独说。其一，策略层只能读引擎给定的数据快照，不能自己取未来数据——look-ahead 一旦混进策略分支，引擎层的常规检查发现不了。其二，每次运行的参数与结果必须留痕：同一个策略在不同参数下跑一百次，账本里就是一百条记录，第~\ref{sec:mtflow}~节的试验族就以这份账本为准；只记最优一次，事后无从还原搜索代价。

\subsection{规则一：期货时序趋势跟踪}

\begin{quote}\ttfamily\small
每个交易日 $t$：\\[2pt]
\quad 若 $t$ 是重平衡日（距上次重平衡满 5 个交易日）：\\[2pt]
\qquad 对候选池中每个品种 $i$：\\
\qquad\qquad $s_i \gets \sum_{k=0}^{L-1}\ln(1+r_{i,t-k})$ \hfill \# 过去 $L$ 日累计对数收益\\
\qquad\qquad $v_i \gets \mathrm{std}(r_i,\ 60\ \text{日})\times\sqrt{252}$ \hfill \# 已实现波动年化\\
\qquad\qquad 若 $|s_i|/(\mathrm{std}(r_i)\times\sqrt L)\le 1.0$：\\
\qquad\qquad\qquad 目标权重 $w_i \gets 0$ \hfill \# 信号强度不足，弃权\\
\qquad\qquad 否则：\\
\qquad\qquad\qquad $w_i \gets \mathrm{sign}(s_i)\times\min(10\%/v_i,\ 3.0)$\\[2pt]
\qquad $w \gets w/\sum_j|w_j|$ \hfill \# 归一化，总敞口为 1\\[2pt]
\qquad 若 $\big(R_{\text{mkt}}^{(60)}<-5\%\big)$ 且 $\big(\sigma_{\text{mkt}}^{(60)}>\text{历史 80 分位}\big)$：\\
\qquad\qquad $w \gets 0.5\times w$ \hfill \# 恐慌状态降敞口\\[2pt]
\qquad 若 $\bar\rho_t>$ 历史 90 分位：\\
\qquad\qquad $w \gets 0.7\times w$ \hfill \# 拥挤度预警\\[2pt]
\qquad 对每个 $i$：\\
\qquad\qquad 若 $|w_i-w_i^{\text{current}}|<\theta_{\text{band},i}$：\\
\qquad\qquad\qquad $w_i \gets w_i^{\text{current}}$ \hfill \# 免交易带\\[2pt]
\qquad 提交调仓指令：\\
\qquad\qquad 对每个 $i$，若当日涨跌停或停牌：顺延至下一可交易日\\
\qquad\qquad 否则按 $w_i$ 下单，收取单边 2 基点成本\\[2pt]
\quad 更新持仓漂移：\\
\qquad 对每个 $i$：$w_i \gets w_i\times(1+r_{i,t})/\big(1+\sum_j w_j r_{j,t}\big)$\\[2pt]
\quad 检查防线：\\
\qquad 若组合回撤 $\ge15\%$：全部仓位 $\times0.5$\\
\qquad 若组合回撤 $\ge25\%$：全部平仓，暂停 20 个交易日\\
\qquad 对每个 $i$，若浮亏 $\ge2.5\times\text{ATR}_{20}(i)$：该品种平仓
\end{quote}

\subsection{规则二：A 股截面反转}

\begin{quote}\ttfamily\small
每 21 个交易日（调仓日）：\\[2pt]
\quad 候选池 $\gets$ 全部 A 股\\[2pt]
\quad 过滤：\\

\qquad 剔除 ST 与 *ST\\
\qquad 剔除上市不足 250 个交易日\\
\qquad 剔除近 20 日日均成交额 $<5000$ 万元\\
\qquad 标记当日涨跌停与停牌股票（不参与本次调仓）\\[2pt]
\quad 对候选池每个 $i$：\\
\qquad $s_i \gets \sum_{k=0}^{L-1}\ln(1+r_{i,k})$ \hfill \# 过去 $L$ 日累计对数收益（首选 $250$，备选 $60$）\\[2pt]
\quad 横截面回归取残差：\\
\qquad $s_i \gets \alpha+\sum_g\beta_g\mathbf 1_{\{i\in g\}}+\gamma\ln(\mathrm{MV}_i)+\tilde s_i$\\
\qquad 用残差 $\tilde s_i$ 排序 \hfill \# 行业 + 市值中性化\\[2pt]
\quad 多头集合 $\gets$ 残差最低的 20\%，按实盘约束收窄至 $30\sim50$ 只\\
\quad 目标权重 $\gets$ 多头集合内等权 \hfill \# 版本 A：纯多头\\[2pt]
\quad 计算对冲量：\\
\qquad $\beta_{\text{port}} \gets$ 组合对目标指数（IC/IF/IM）的 60 日滚动 beta\\
\qquad 空头名义 $\gets \beta_{\text{port}}\times$ 多头市值 \hfill \# 指数对冲\\[2pt]
\quad 若组合回撤 $\ge20\%$：目标敞口 $\times0.5$\\
\quad 若组合回撤 $\ge30\%$：全部平仓，暂停 40 个交易日\\[2pt]
\quad 执行调仓（分 $2\sim3$ 个交易日完成，收取单边 15 基点成本）\\[2pt]
每个交易日：\\
\quad 对每只持仓股票：\\
\qquad 若相对建仓加权成本跌幅 $\ge15\%$（高波动期 $\ge10\%$）：该票止损\\
\qquad 若创 120 日新低且跌破 60 日最低价：该票止损\\
\qquad 【硬约束】不得因跌幅扩大而加仓；仅当已从低点反弹 $\ge5\%$ 方可考虑补仓
\end{quote}

\subsection{规则三：期货截面统计套利}

\begin{quote}\ttfamily\small
每 21 个交易日（滚动重建）：\\[2pt]
\quad 按产业链分组候选品种\\[2pt]
\quad 对每个产业链内候选对 $(A,B)$：\\
\qquad 取形成期 60 日数据\\
\qquad $\beta \gets \mathrm{OLS}\big(\ln P^A \sim \ln P^B\big)$\\
\qquad $e \gets \ln P^A-\beta\ln P^B$\\
\qquad 若 $\mathrm{ADF}(e)\ge$ 5\% 临界值：跳过该对 \hfill \# 协整资格检验\\
\qquad $\phi \gets$ 对 $e$ 拟合 AR(1) 的系数\\
\qquad 若 $\phi\notin(0,1)$：跳过该对\\
\qquad $\theta \gets -\ln(\phi)/\Delta$；\ $H \gets \ln 2/\theta$\\
\qquad 若 $H\notin[5,60]$ 交易日：跳过该对 \hfill \# 半衰期筛选\\
\qquad 记录 $(A,B,\alpha,\beta,H,\mu_e,\sigma_e)$ 进入活跃对列表\\[2pt]
每个交易日：\\[2pt]
\quad 对每个活跃对：\\
\qquad 用 Kalman 滤波更新 $(\alpha_t,\beta_t)$ \hfill \# 在线对冲比\\
\qquad $e_t \gets \ln P^A_t-\alpha_t-\beta_t\ln P^B_t$\\
\qquad $\mu_t \gets$ 价差滚动 60 日均值；$\sigma_t \gets$ 滚动 60 日标准差\\
\qquad $z_t \gets (e_t-\mu_t)/\sigma_t$\\[2pt]
\qquad 若无持仓 且 $|z_t|>2.0$：\\
\qquad\qquad 开仓：$z_t>2$ 时空 $A$ 多 $B$；$z_t<-2$ 时多 $A$ 空 $B$\\
\qquad\qquad 仓位按价差波动率倒数归一\\
\qquad 若有持仓 且 $|z_t|<0.5$：平仓 \hfill \# 收敛了结\\
\qquad 若有持仓 且 $|z_t|>3.5$：强制止损 \hfill \# 不收敛防线\\
\qquad 若有持仓 且持仓天数 $>3H$：平仓 \hfill \# 时间止损\\[2pt]
\quad 组合层面：\\
\qquad 同一产业链持仓对数 $\le$ 总对数的一半\\
\qquad 若组合回撤 $\ge20\%$：敞口 $\times0.5$\\
\qquad 若组合回撤 $\ge30\%$：全部平仓，暂停 40 个交易日\\
\qquad 保证金冗余 $\ge$ 名义敞口的 40\%\\[2pt]
\quad 成本参数：\\
\qquad 单边 2 基点；若实际成本升至 $c'$，\\
\qquad 按 $b\propto c^{1/3}$ 把开仓阈值由 2.0 调整为 $2.0\times(c'/c)^{1/3}$
\end{quote}

\subsection{三套伪代码的共同结构}

三段并排看，骨架完全相同：\textbf{信号 $\to$ 归一化/中性化 $\to$ 缩放决策 $\to$ 门槛过滤 $\to$ 执行 $\to$ 防线}。差异只在三处——规则一有波动率缩放与免交易带、门槛是信号强度；规则二、三不做缩放，规则三的门槛是协整与半衰期、规则二没有独立门槛；防线上规则一是事前的状态过滤与拥挤监控，规则二、三是事中的止损与资金冗余。这三处正是第~\ref{sec:compare}~节“四行差异”里的三行，其余（成本处理、涨跌停顺延、调仓节奏、评估口径）三段完全一致——这是论文“唯一区别是符号，其余差异全在执行层”在规则层面的直接呈现。

\section{规则三的协整配对筛选算法}\label{sec:pairalgo}

本节给出规则三配对筛选的完整算法步骤。筛选分七步，每步都给出判据、实现要点与拒绝率实测。实现在 \texttt{code/rule3\_pairs.py}，可直接运行（\texttt{python rule3\_pairs.py --screen}）。

\subsection{算法的整体结构}

筛选的目的是从产业链内的两两组合中，找出「价差确实平稳、且回复速度落在可交易区间」的那些对。七步的顺序不可调换：先分组降维、再初筛降噪、再严格检验、最后估计参数。

\begin{quote}\ttfamily\small
S1 产业链分组 $\to$ 生成候选对\\
S2 相关性初筛（去伪配对）\\
S3 Engle--Granger 协整检验（OLS + ADF）\\
S4 OU 参数估计（AR(1) $\to$ $\theta$、半衰期）\\
S5 半衰期区间筛选 $[5,60]$ 交易日\\
S6 滚动重估（形成期 60 日 / 交易期 21 日）\\
S7 生成可交易对表（对冲比率、均值、标准差、半衰期）
\end{quote}

\subsection{S1：产业链分组与候选对生成}

\textbf{判据}：只在同一产业链内配对。跨产业链的配对数不足信——两条序列之间缺少经济意义上的长期均衡关系，协整容易被伪回归误判。

\textbf{实现}：按表~\ref{tab:chains}~的五条产业链展开。设第 $c$ 条链有 $n_c$ 个品种，则该链产生 $\binom{n_c}{2}$ 对；去重后汇总。

\textbf{实测}：五条产业链共 118 个候选对。

\subsection{S2：相关性初筛}

\textbf{判据}：形成期两序列对数价格的相关性 $\rho\ge0.6$。

\textbf{理由}：这一步不是统计检验，而是计算量与稳健性的过滤。相关性过低的配对在后续 ADF 检验中几乎必然被拒，先剔除以节省计算；更重要的是，低相关配对即便偶然通过协整检验，其价差波动也极大，交易风险不可控。

\textbf{实现}：
\begin{equation}
\rho_{AB}=\mathrm{Corr}\bigl(\ln P^A_t,\ \ln P^B_t\bigr),\qquad t\in\text{形成期}.
\end{equation}

\textbf{实测}：16992 次检验中被此步拒绝 8760 次（$51.6\%$），是最主要的过滤环节。

\subsection{S3：Engle--Granger 协整检验}

\textbf{判据}：对残差做 ADF 检验，$p<0.05$。这是本算法唯一的「资格检验」，其依据是反转类策略的条件性特征（论文对偶七）：\emph{必须先证明价差会回归，再交易}。

\textbf{实现}（两步法\cite{englegranger1987}）：

\begin{enumerate}[leftmargin=2.2em]
  \item 对形成期做 OLS：$\ln P^A_t=\alpha+\beta\ln P^B_t+e_t$，取 $\hat\beta$；
  \item 对残差 $e_t$ 做 ADF 检验（含常数项、滞后阶数由 AIC 自动选择\cite{dickeyfuller1979}），
        统计量需小于 $5\%$ 临界值；等价地 $p<0.05$。
\end{enumerate}

\textbf{实现要点}：OLS 用最小二乘直接求解（\texttt{lstsq}），ADF 用统计库的 \texttt{adfuller}。残差必须做去均值处理，否则 ADF 的常数项设定与协整回归的截距重复。

\textbf{已知局限}：Engle--Granger 两步法在样本量较小时检验功效偏低（论文附录 A.1 讨论的 $\hat\phi$ 下偏问题同源）。若形成期只有 60 个观测，ADF 的临界值并非标准分布，严格做法应使用 Engle--Granger 专用的响应面临界值（比标准 ADF 临界值更负）。本实现在此采用标准 ADF 临界值，属\emph{偏宽松}的近似，会略微高估通过率。

\textbf{实测}：被此步拒绝 4942 次（$29.1\%$）。通过者的平均 ADF $p$ 值为 $0.0203$。

\subsection{S4：OU 参数估计}

\textbf{判据}：残差可拟合为稳定 AR(1)，即 $0<\phi<1$。

\textbf{实现}：对残差序列拟合
\begin{equation}
e_{t+1}=c+\phi e_t+u_t ,
\end{equation}
由 OLS 得到 $\hat\phi$（同样可用极大似然，二者在正态假设下渐近等价）。再换算为连续时间参数：
\begin{equation}
\theta=-\frac{\ln\hat\phi}{\Delta},\qquad
\mu=\frac{c}{1-\phi},\qquad
H=\frac{\ln 2}{\theta}\quad\text{（半衰期）}.
\end{equation}

\textbf{为何必须检查 $0<\phi<1$}：$\phi\ge1$ 意味着残差非平稳（价差不回归），$\phi\le0$ 意味着残差是高频振荡而非缓慢回归，两者都使半衰期失去意义。论文附录指出 $\hat\phi$ 在小样本下有向下偏误，因此短形成期估出的半衰期偏短——这一点在 S5 的区间下界设置中有直接后果（见下）。

\textbf{实测}：仅 1 次被此步拒绝，说明在通过 ADF 的基础上 $\hat\phi$ 几乎总落在 $(0,1)$ 内。

\subsection{S5：半衰期区间筛选}

\textbf{判据}：$H\in[5,\,60]$ 交易日。

\textbf{两侧边界的含义}：
\begin{itemize}[leftmargin=2.2em]
  \item \textbf{下界 $H\ge5$}：半衰期太短意味着价差回复极快，交易频率上升、成本侵蚀收益。这与论文铁律二一致——本文实测的 A 股短窗口反转（$L=5$）毛利夏普 $1.08$ 被成本压至 $-0.35$，同类机制在此适用。
  \item \textbf{上界 $H\le60$}：半衰期太长意味着占用保证金数月才收敛一半，资金效率过低，且期间发生结构性断裂（如政策冲击改变均衡关系）的概率上升。
\end{itemize}

\textbf{实测}：被此步拒绝 1811 次（$10.7\%$）；最终通过 75 条。\textbf{通过者的半衰期中位数仅 5.9 交易日}，紧贴下界——这说明中国商品期货中「通过协整检验」的配对，其价差大多回复极快。这一发现对规则三有实质影响：它意味着通过筛选的配对天然是高换手的，成本敏感度高，与论文铁律二形成张力。

\subsection{S6：滚动重估}

\textbf{判据}：形成期 60 交易日、交易期 21 交易日，滚动推进，两段严格不重叠。

\textbf{理由}：若配对是在包含交易期的数据上选出的，价差的平稳性会被机械地保证，构成 look-ahead。这一「形成期与交易期严格分离」的范式来自 Gatev 等(2006) 的 $6+6$ 个月滚动设计\cite{gatev2006}。

\textbf{实现要点}（本实现的关键正确性问题）：交易期的标准化必须用\emph{滚动}均值与标准差，而非形成期常数。原因很直接——形成期的价差波动无法代表交易期的波动；若用形成期 $\sigma$ 计算 $z$，交易期一旦波动放大，$|z|$ 会虚高，导致「开仓即触止损」。本文首版实现即犯此错，成交笔数虚增、结果不可用；改为「形成期 $+$ 已过去交易期」的滚动 60 日均值与标准差后恢复正常。

\textbf{实测}：形成期数 144 期，累计 16992 次配对检验，通过 75 条。

\subsection{S7：生成可交易对表与阈值交易}

\textbf{产出字段}：形成期末日期、产业链、品种 A/B、相关性、ADF $p$、对冲比率 $\beta$、半衰期 $H$、价差均值 $\mu$、价差标准差 $\sigma$。

\textbf{对冲比率}。静态版用形成期 OLS 的 $\hat\beta$；在线版用 Kalman 滤波滚动估计 $\beta_t$\cite{elliott2005}。Kalman 的状态空间形式为
\begin{equation}
\begin{aligned}
\ln P^A_t&=\alpha_t+\beta_t\ln P^B_t+\varepsilon_t, &&\varepsilon_t\sim\mathcal N(0,\sigma_\varepsilon^2),\\
\alpha_{t+1}&=\alpha_t+\eta^{\alpha}_t,\quad \beta_{t+1}=\beta_t+\eta^{\beta}_t, &&\eta\sim\mathcal N(0,\Sigma_\eta).
\end{aligned}
\end{equation}

\textbf{必须使用一步前瞻的 $\beta_t$}。这是本实现的第二个关键正确性问题：若在时刻 $t$ 先用第 $t$ 期观测更新 $\beta$，再用更新后的 $\beta$ 计算价差 $e_t$，则 $e_t$ 被当期价格拟合，残差呈现\emph{人为的}均值回归，回测会给出虚高的夏普。本文首版即因此得到夏普 $3.51$ 的荒谬结果；改为「先记录前瞻 $\beta$、再用当期观测更新」后，夏普降至 $1.09$，与成交笔数（21 笔）相符。

\textbf{阈值规则}（把连续带宽离散化为事件规则）：

\begin{table}[htbp]
\centering\footnotesize
\caption{规则三的进出场阈值}\label{tab:pairband}
\begin{tabular}{@{}>{\raggedright\arraybackslash}p{3.0cm}>{\raggedright\arraybackslash}p{5.4cm}>{\raggedright\arraybackslash}p{5.5cm}@{}}
\toprule
\textbf{条件} & \textbf{动作} & \textbf{依据} \\
\midrule
$|z_t|>2.0$ 且空仓 & 开仓（$z>2$ 空 A 多 B；$z<-2$ 反之） & s-score 范式\cite{avellaneda2010} \\
$|z_t|<0.5$ 且持仓 & 平仓（收敛了结） & 同上 \\
$|z_t|>3.5$ 且持仓 & 强制止损（不收敛防线） & \cite{leungli2015} \\
持仓 $>3H$ 且未收敛 & 平仓（时间止损） & 资金效率 \\
实际成本升至 $c'$ & 开仓阈值改为 $2.0\times(c'/c)^{1/3}$ & 带宽随成本立方根\cite{martin2011} \\
\bottomrule
\end{tabular}
\end{table}

\textbf{风险归一}。每对按风险预算归一：设 1 倍标准差的价差变动对应 \texttt{risk\_per\_pair} 的资金敞口，则日损益
\begin{equation}
\Delta\text{PnL}_t=\text{direction}\times\Delta z_t\times\text{risk\_per\_pair}.
\end{equation}
总风险预算按同时持仓上限均分（本实现：总风险 $30\%$，上限 20 对，故每对 $1.5\%$）。若不归一，多对敞口直接相加会使净值波动失控（本文首版曾得到 $-100.96\%$ 的「最大回撤」，即净值归零）。

\textbf{实测结果}：成交 21 笔，平均持仓 11.6 天，有仓日的平均同时持仓对数 1.2，年化夏普 $1.09$、年化 $4.31\%$、最大回撤 $-4.06\%$。

\textbf{必须说明的局限}：21 笔交易在 12 年样本上不足以支撑统计推断。报告中 $t=3.99$ 的 Newey--West 统计量是按 \emph{3039 个交易日}计算的，而策略在绝大多数交易日空仓，该 $t$ 值严重高估显著性；正确的检验应基于\emph{交易笔数}（$n=21$），其功效极低。因此本节的规则三结果只能读作「筛选流程可用」，不能读作「策略有效」。这也正是下一节多重检验流程要处理的问题。

\subsection{筛选拒绝率的诊断价值}

各步拒绝率的分布本身就是信息：

\begin{table}[htbp]
\centering\footnotesize
\caption{配对筛选各步拒绝率实测（16992 次检验）}\label{tab:reject}
\begin{tabular}{@{}>{\raggedright\arraybackslash}p{2.4cm}>{\raggedright\arraybackslash}p{3.4cm}>{\raggedright\arraybackslash}p{2.4cm}>{\raggedright\arraybackslash}p{6.0cm}@{}}
\toprule
\textbf{步骤} & \textbf{判据} & \textbf{拒绝次数（占比）} & \textbf{解读} \\
\midrule
S2 & $\rho\ge0.6$ & 8760（$51.6\%$） & 产业链内仍有大量品种对走势独立，先验均衡关系薄弱 \\
S3 & ADF $p<0.05$ & 4942（$29.1\%$） & 高相关不等于协整，这一步是真正的资格门槛 \\
S5 & $H\in[5,60]$ & 1811（$10.7\%$） & 通过协整者多为极快回复（中位数 5.9 日），被下界拦下一部分 \\
S4 & $0<\phi<1$ & 1（$0.0\%$） & ADF 通过后残差基本稳定 \\
\midrule
合计通过 & — & 75（$0.44\%$） & 每 227 次检验仅有 1 条可交易对 \\
\bottomrule
\end{tabular}
\end{table}

「每 227 次检验仅 1 条通过」这一比例，正好印证论文对偶七的判断：反转类型是\emph{条件性}策略，它先要在一大批候选中筛掉绝大多数，剩下的才有资格进入交易设计。这也解释了为什么反转策略的有效样本基础天然比趋势策略薄。

\section{上线前的多重检验校正流程}\label{sec:mtflow}

论文在审计章明确指出：自建实证未做多重检验校正，$48$（实际为 $60$）个参数组合的网格扫描结果不能读作统计显著。本节把这句警告落实为一个可执行的流程，并对 $60$ 个组合实跑。

\subsection{流程设计：四道闸门}

流程的目的不是「证明策略有效」，而是\emph{找出那些在数据窥探下依然站得住的结果}。四道闸门的判据不同，必须同时通过：

\begin{enumerate}[leftmargin=2.2em]
  \item \textbf{闸门一：White Reality Check}——对「所有策略的期望收益均 $\le0$」这一零假设做自举检验\cite{white2000}。它回答：$60$ 个组合里最好的那个，是否只是搜索空间的产物？
  \item \textbf{闸门二：Hansen SPA}——Reality Check 的 studentized 改进版，对表现差的策略做再中心化，避免劣质策略拖累检验\cite{hansen2005}。
  \item \textbf{闸门二附：Romano--Wolf 逐步多重检验}——给出\emph{每一个}假设的逐步校正 $p$ 值，在控制族错误率（FWER）的同时比单步法保留更多功效\cite{romanowolf2005}。
  \item \textbf{闸门三：紧缩夏普比率（DSR）}——把「试验次数」直接计入显著性：把最优夏普与「$N$ 次独立试验下的期望最大夏普」比较\cite{bailey2014}。
  \item \textbf{闸门四：回测过拟合概率（PBO/CSCV）}——用组合对称交叉验证估计「样本内最优在样本外落到中位数以下」的概率\cite{bailey2016}。
\end{enumerate}

四道闸门的关系需要说清：闸门一、二检验\emph{最大值}的显著性，闸门三检验\emph{幅度}是否超过搜索带来的期望膨胀，闸门四检验\emph{选择过程}是否稳健。它们测的是不同侧面，因此可能出现分歧——本文的实测正好出现了分歧。

\subsection{方法细节}

\subsubsection*{平稳自举}

Reality Check 与 SPA 都需要对时间序列重采样。金融收益存在波动率聚集，简单 i.i.d. 自举会破坏这一结构、低估统计量的方差。本文采用 Politis--Romano 平稳自举：以几何分布的块长抽样，平均块长取 20 个交易日，自举次数 2000。

\subsubsection*{Reality Check 统计量}

非 studentized 版本：
\begin{equation}
V=\max_k\ \sqrt{T}\,\bar f_k ,
\end{equation}
其中 $\bar f_k$ 为第 $k$ 个组合的日均收益。自举时以 $\bar f_k$ 为中心再抽样，得到 $V^*_b=\max_k\sqrt T\,(\bar f^{*}_{k,b}-\bar f_k)$，$p$ 值为 $\Pr(V^*_b\ge V)$。

studentized 版本把 $\bar f_k$ 换成 $t_k=\bar f_k/\hat\sigma_k$（$\hat\sigma_k$ 为 Newey--West 标准误），量纲无关，因此不会因为某些组合的波动极大而主导最大值。

\textbf{为何两个版本都要报}：非 studentized 版本对高波动组合极其敏感——本文的 $60$ 个组合横跨期货与 A 股两类资产，波动相差数倍，非 studentized 统计量几乎被波动最大的组合支配（实测统计量仅 $0.02$，即最大值来自一个低波动组合）。

\subsubsection*{Romano--Wolf 逐步检验}

单步的 Reality Check 只回答「最强的那个是否显著」；实务上更常需要「哪些组合能被留下」。Romano--Wolf(2005) 给出逐步算法\cite{romanowolf2005}：

\begin{enumerate}[leftmargin=2.2em]
  \item 把 $t$ 统计量由大到小排序，得假设序列 $(1),(2),\dots,(N)$；
  \item 用\emph{同一次}自举抽样计算再中心化的 studentized 统计量矩阵 $t^*_{b,k}$；
  \item 第 $j$ 步只在剩余假设集合 $\{(j),(j+1),\dots,(N)\}$ 上取
        $\max_{k\ge j}t^*_{b,k}$ 作为临界分布，与 $\max_{k\ge j}t_{(k)}$ 比较得逐步 $p$ 值；
  \item 一旦某步 $p>\alpha$ 即停止，其后假设一律不拒绝。
\end{enumerate}

第 3 步的「逐步缩小比较集」是功效来源：越靠后的假设面对的临界值越松。第 2 步的「共用同一次自举」是它与「逐条单独检验再加 Bonferroni」的本质区别——后者会因忽略假设间相关性而过度保守。

\textbf{一个可用的实现校验}：RW 第 1 步的 $p$ 值应当\emph{恰好等于} studentized Reality Check 的 $p$ 值（两者在第 1 步定义相同）。本文实测两者均为 $0.0455$，这构成对实现的内部一致性检验。

\subsubsection*{Hansen SPA}

SPA 的关键是再中心化：对 $\hat t_k$ 低于阈值 $-\sqrt{2\log\log T}$ 的「劣质策略」，自举时用 $0$ 而非 $\bar f_k$ 作中心，从而使其不再贡献正的偏差。统计量取 $\max(0,\max_k t_k)$。

\subsubsection*{紧缩夏普比率}

令 $\widehat{SR}$ 为最优组合的（日频）夏普，$T$ 为样本长度，$\hat\gamma_3$、$\hat\gamma_4$ 为其偏度与峰度，$\hat V$ 为全部 $N$ 次试验夏普的横截面方差。期望最大夏普为
\begin{equation}
SR_0=\sqrt{\hat V}\left[(1-\gamma)\,\Phi^{-1}\!\left(1-\frac1N\right)
+\gamma\,\Phi^{-1}\!\left(1-\frac{1}{N e}\right)\right],
\qquad \gamma=0.5772\ldots ,
\end{equation}
紧缩夏普比率为
\begin{equation}
\text{DSR}=\Phi\!\left(\frac{\left(\widehat{SR}-SR_0\right)\sqrt{T-1}}
{\sqrt{1-\hat\gamma_3\widehat{SR}+\frac{\hat\gamma_4-1}{4}\widehat{SR}^2}}\right).
\end{equation}
$\text{DSR}>0.95$ 才可认为最优夏普在 $5\%$ 水平上超过搜索带来的期望膨胀。

\textbf{口径一致性的陷阱}：$SR_0$ 与 $\widehat{SR}$ 必须同为日频（或同为年化）。本文首版把年化的试验方差与日频的 $\widehat{SR}$ 混用，得到 $SR_0=1.605$（日频！）的荒谬值，使 DSR 恒为 $0$。修正口径后 $SR_0=0.1011$，结果才有意义。这类单位错误在多重检验实现中极易发生，且因为「结论都是拒绝」而不会被察觉，需要专门校验。

\subsubsection*{PBO/CSCV}

把 $T$ 个交易日切成 $S=10$ 个等长区块，枚举所有 $\binom{10}{5}=252$ 种「5 块作样本内、5 块作样本外」的切分。对每种切分：在样本内选出夏普最高的组合 $n^*$，计算其在样本外的夏普排名 $\omega\in(0,1)$，取 $\lambda=\ln\frac{\omega}{1-\omega}$。PBO 为 $\lambda<0$ 的比例，即样本内最优在样本外落到中位数以下的概率。$\text{PBO}>0.5$ 表示选择过程不比抛硬币更好。

\subsection{实测结果}

对 $60$ 个组合（$4$ 组口径 $\times$ 两个方向 $\times$ 六个窗口）实跑：

\begin{table}[htbp]
\centering\footnotesize
\caption{多重检验校正实跑结果（$60$ 个组合，$3039$ 个交易日）}\label{tab:mtres}
\begin{tabular}{@{}>{\raggedright\arraybackslash}p{4.1cm}>{\raggedright\arraybackslash}p{3.2cm}>{\raggedright\arraybackslash}p{1.9cm}>{\raggedright\arraybackslash}p{5.0cm}@{}}
\toprule
\textbf{方法} & \textbf{统计量} & \textbf{结果} & \textbf{判读} \\
\midrule
White RC（非 studentized） & $V=0.02$ & $p=0.523$ & 不拒绝；但该统计量被波动结构支配，参考价值低 \\
White RC（studentized） & $V=3.18$ & $p=0.046$ & 边际拒绝（$5\%$ 线上） \\
Hansen SPA & $V=3.18$ & $p=0.038$ & 拒绝；10/60 被判为劣质策略 \\
Romano--Wolf（FWER $5\%$） & 逐步 & 拒绝 $1/60$ & 仅 \texttt{astk\_XS\_rev\_L60} 存活；第 2 步 $p=0.146$ 即停止 \\
紧缩夏普比率（DSR） & $SR_0=0.1011$ & $0.016$ & \textbf{不通过}：最优日夏普 $0.0625$ 低于 $60$ 次试验的期望最大值 \\
PBO / CSCV（$S=10$） & $252$ 种切分 & $0.111$ & 通过：样本内最优在样本外平均排在第 $84.2$ 百分位 \\
\bottomrule
\end{tabular}
\end{table}

Romano--Wolf 的结果给出一张更清晰的图景：\textbf{$60$ 个组合中只有 $1$ 个能在族错误率 $5\%$ 下存活}，且第二名（\texttt{L=120}）的逐步 $p$ 值立刻跳到 $0.146$——优势并不连续。这与「哪些参数值得留下」的实务判断直接衔接：可留下的只有一个点，而一个孤立的点恰恰是稳健性最差的情形。

最优组合为 \texttt{A股个股截面反转 L=60}（年化夏普 $0.99$），与论文表 $6$ 一致——这同时验证了本文网格与论文口径的可复现性。

\subsection{分歧的解读：为什么 DSR 拒绝而 SPA 与 PBO 通过}

四项结果并不一致，这不是实现错误，而是它们测的东西不同。这个分歧本身值得展开。

\textbf{SPA 与 studentized RC 通过了}，因为它们检验的是「最大值是否显著为正」。实测 $t=3.18$ 的自举分布下，$p$ 值确实落在 $5\%$ 附近。也就是说：\emph{在 $60$ 个组合中，确实存在一个统计上强于零的组合}。

\textbf{DSR 拒绝，因为它引入了幅度基准}。$SR_0=0.1011$ 是「从 $N(0,\hat V)$ 中抽 $60$ 次取最大」的期望值——由于 $60$ 个试验的夏普横截面方差较大，这个期望最大值本身就很高（日频 $0.10$，年化约 $1.6$）。最优组合的日夏普 $0.0625$ 反而低于它。换句话说：\emph{最优组合的夏普绝对值，与「纯粹靠搜索挑出来的最大值」没有区别}。

\textbf{PBO 通过，因为它测的是排序而非幅度}。样本内最优在样本外平均排在第 $84.2$ 百分位，说明选择过程是有信息的（不是随机的）。但它不回答「这个排名对应的收益是否值得交易」。

三者的合并结论应当是：\textbf{A 股 $L=60$ 反转确实是一个「相对更好」的组合，但它的夏普幅度不足以排除「它只是搜索空间的极值」这一解释}。这正好对应论文结论置信度一章给该口径的评级——实证直接支持其\emph{相对}优势，但绝对水平不可承诺。

\subsection{流程的落地形式}

把上述四道闸门写成上线前的固定步骤：

\begin{quote}\ttfamily\small
步骤 1：登记搜索空间\\
\quad 记录本次策略研发的全部试验次数 N（含失败与废弃的参数）\\
\quad 未登记的搜索空间会使后续校正失效\\
步骤 2：计算网格日收益矩阵（T × N）\\
步骤 3：闸门一、二（自举类：RC / SPA / Romano--Wolf）\\
\quad 若三者均不拒绝：停止，判定为数据窥探产物\\
步骤 4：闸门三（DSR）\\
\quad 若 DSR < 0.95：不得以「统计显著」表述，降级为「候选」\\
步骤 5：闸门四（PBO）\\
\quad 若 PBO > 0.5：判定选择过程无效，需重新设计信号而非调参\\
步骤 6：对通过者施加样本外检验（滚动 walk-forward）与\\
\quad 可交易性过滤（涨跌停、停牌、参与率约束）\\
步骤 7：归档：N、各条 p 值、DSR、PBO、样本外结果，一并写入策略档案
\end{quote}

\textbf{步骤 1 是最容易被跳过、也最关键的一步}。多重检验校正的全部依据是「试验次数 $N$」；如果研发过程中试过的参数、改过的信号、调过的止损没有全部登记，$N$ 会被系统性低估，校正随之失效。本文的 $N=60$ 是完整登记的网格，因此 DSR 的判定可信；若实际研发中还有未登记的尝试，真实的 $SR_0$ 会更高，结论只会更保守。

\subsection{总入口：一次跑完七步并出总判决}

上述七步此前分散在三个脚本里（\texttt{multipletest.py} 四道闸门、\texttt{walkforward.py} 组合层面的样本外、\texttt{rule\_wf.py} 规则自身的样本外）。现已收敛为唯一入口 \texttt{code/go\_live\_check.py}，跑一次输出逐项结论与总判决：

\begin{quote}\ttfamily\small
python go\_live\_check.py \hfill \# 完整检查（自举 2000 次，约 2 分钟）\\
python go\_live\_check.py --fast \hfill \# 快速检查（自举 300 次，约 10 秒）\\
python go\_live\_check.py --with-rule-wf \hfill \# 额外跑规则自身参数的样本外检验
\end{quote}

\textbf{退出码约定}：$0$ = 七步全部通过；$1$ = 存在未通过项。这一约定使它可以\emph{直接挂在 CI 里做阻断}——研究流程产出的新参数组合若过不了闸门，合并请求就无法通过。

本文对当前网格（$N=60$）实跑的结果：

\begin{table}[htbp]
\centering\footnotesize
\caption{上线检查单实跑结果（$N=60$，自举 300 次）}\label{tab:golive}
\begin{tabular}{@{}>{\raggedright\arraybackslash}p{2.4cm}>{\raggedright\arraybackslash}p{1.4cm}>{\raggedright\arraybackslash}p{9.6cm}@{}}
\toprule
\textbf{步骤} & \textbf{结果} & \textbf{依据} \\
\midrule
1 搜索空间登记 & ✓ & $N=60$、$T=3039$ \\
2 网格构建 & ✓ & 形状与缺失率检查通过（缺失 $1.37\%$） \\
3 自举类检验 & ✓ & studentized RC $p=0.027$、SPA $p=0.023$、RW 拒绝 $1/60$ \\
4 紧缩夏普比率 & \textbf{✗} & DSR $=0.016$，远低于 $0.95$ \\
5 过拟合概率 & ✓ & PBO $=0.111$ \\
6 样本外检验 & \textbf{✗} & IS $\to$ OOS 相对衰减 $54.7\%$，超过 $50\%$ 判据 \\
7 可交易性核查 & \textbf{✗} & 涨跌停与停牌过滤尚未接入回测引擎（引擎来自论文口径） \\
\midrule
\multicolumn{3}{@{}l}{\textbf{总判决：$3$ 项未通过，不得直接上线 —— DSR、样本外、可交易性}} \\
\bottomrule
\end{tabular}
\end{table}

这张表比任何单一指标都更能说明问题：\textbf{「最强的组合在自举下确实显著」（步骤 3 通过），但它同时过不了幅度基准（步骤 4）、样本外外推（步骤 6）与可执行性（步骤 7）}。这正是第~\ref{sec:verify}~节反复出现的形态——单看一项容易得出乐观结论，放到完整闸门序列里就会露馅。

步骤 7 的未通过是\textbf{刻意保留的诚实结果}：本文网格由论文口径的引擎生成，该引擎未做涨跌停与停牌的成交剔除。要让它通过，必须把共享层的 \texttt{C.tradable\_mask} 接进成交判定（涨跌停当日顺延、停牌不调仓），并把 \texttt{go\_live\_check.py} 顶部的 \texttt{TRADABILITY\_APPLIED} 置为 \texttt{True} 后重跑。这个开关刻意做成人工确认，以免「声称做了过滤」而实际没做。

\subsection{回归测试}

四道闸门实现在 \texttt{code/multipletest.py}，运行 \texttt{python multipletest.py} 即可复现本节全部数字（$N=60$，自举 $2000$ 次，随机种子固定为 $20261008$）。框架另附 \texttt{code/tests/test\_rules.py}（$23$ 项，全部通过，运行约 $3$ 秒），覆盖四类：

\begin{itemize}[leftmargin=2.2em]
  \item \textbf{共享层数值正确性}：涨跌停分板块判定、指标与手算值一致、成本门槛边界；
  \item \textbf{策略对称性}（G1 门禁）：权重取负后毛收益是否\emph{严格}互为相反数、换手是否完全相同。
        测试同时揭示了一个此前未被写明的边界——\emph{含权重漂移时对称性只是近似}（漂移分母 $1\pm r$ 不同），
        只有关闭漂移才是精确的；
  \item \textbf{统计工具性质}：EG 临界值在 $T=60$ 处等于 $-3.4396$、比标准 ADF 更负、随样本长度单调；OU 参数能回收已知半衰期；Kalman 返回的一步前瞻 $\beta$ 首项等于初值（不吸收当期观测）；
  \item \textbf{多重检验思想校验}：纯噪声下 RC 不拒绝、注入强信号后拒绝；SPA 不会把劣质策略「再中心化」成显著；RW 在噪声下零拒绝、在强信号下捕获该列；DSR 在 $N=1$ 时退化为标准正态检验且随试验数增加而下降；PBO 在「全样本一致的真信号」下偏低、在「仅样本内有效」下偏高。
\end{itemize}

这些测试的价值不在于覆盖率，而在于把三次实跑中暴露过的错误固化为断言：Kalman 的一步前瞻性质、EG 与标准 ADF 的临界值差、$N=1$ 时 DSR 的退化行为，都曾是实际踩过的坑。三处口径问题——$SR_0$ 与 $\widehat{SR}$ 的单位混用、跨资产网格必须报告 studentized 版本、多数交易日空仓的策略不能用日频 $t$ 值——分别在第~\ref{sec:mtflow}~节与第~\ref{sec:pairalgo}~节展开，此处不再重复。

\section{四项验证补充}\label{sec:verify}

第~\ref{sec:mtflow}~节把「样本外检验」写进了流程的第 6 步，但当时并未实跑。本节补齐四项验证：样本外 walk-forward、行业中性化口径对照、规则一完整窗口网格、规则三统计口径修正。

四项验证的结论并不温和——它们各自修正了前文的一处判断。本文如实报告。

\subsection{样本外 walk-forward：IS 夏普 1.26 对 OOS 0.57}

\textbf{设计}。把 $3039$ 个交易日切成滚动折：样本内窗口 $756$ 日（约 3 年）用于从 $60$ 个组合中选出最优，紧随其后的样本外窗口 $252$ 日（约 1 年）用于检验该选择，步长 $252$ 日。IS 与 OOS 严格不重叠，共得 $9$ 折。实现见 \texttt{code/walkforward.py}。

\begin{table}[htbp]
\centering\footnotesize
\caption{walk-forward 逐折结果（$9$ 折）}\label{tab:wf}
\begin{tabular}{@{}>{\raggedright\arraybackslash}p{3.1cm}>{\raggedright\arraybackslash}p{3.1cm}>{\raggedright\arraybackslash}p{3.5cm}>{\raggedright\arraybackslash}p{1.5cm}>{\raggedright\arraybackslash}p{1.6cm}@{}}
\toprule
\textbf{IS 区间} & \textbf{OOS 区间} & \textbf{选中组合} & \textbf{IS 夏普} & \textbf{OOS 夏普} \\
\midrule
2014-04 $\sim$ 2017-05 & 2017-05 $\sim$ 2018-05 & \texttt{astk\_XS\_rev\_L120} & $1.96$ & $-0.55$ \\
2015-04 $\sim$ 2018-05 & 2018-05 $\sim$ 2019-06 & \texttt{astk\_XS\_rev\_L60} & $1.55$ & $+3.03$ \\
2016-04 $\sim$ 2019-06 & 2019-06 $\sim$ 2020-06 & \texttt{fut\_XS\_rev\_L60} & $1.37$ & $+0.91$ \\
2017-05 $\sim$ 2020-06 & 2020-06 $\sim$ 2021-06 & \texttt{astk\_XS\_rev\_L60} & $0.99$ & $-0.05$ \\
2018-05 $\sim$ 2021-06 & 2021-07 $\sim$ 2022-07 & \texttt{aidx\_TS\_rev\_L5} & $0.80$ & $-0.83$ \\
2019-06 $\sim$ 2022-07 & 2022-07 $\sim$ 2023-07 & \texttt{astk\_TS\_rev\_L250} & $0.90$ & $+1.48$ \\
2020-06 $\sim$ 2023-07 & 2023-07 $\sim$ 2024-08 & \texttt{astk\_TS\_rev\_L250} & $1.05$ & $-0.18$ \\
2021-07 $\sim$ 2024-08 & 2024-08 $\sim$ 2025-08 & \texttt{astk\_XS\_rev\_L120} & $1.26$ & $+1.49$ \\
2022-07 $\sim$ 2025-08 & 2025-08 $\sim$ 2026-09 & \texttt{astk\_XS\_rev\_L120} & $1.46$ & $-0.16$ \\
\bottomrule
\end{tabular}
\end{table}

\textbf{核心数字}：

\begin{table}[htbp]
\centering\footnotesize
\caption{walk-forward 汇总指标}\label{tab:wfsum}
\begin{tabular}{@{}>{\raggedright\arraybackslash}p{5.2cm}>{\raggedright\arraybackslash}p{2.4cm}>{\raggedright\arraybackslash}p{6.4cm}@{}}
\toprule
\textbf{指标} & \textbf{数值} & \textbf{含义} \\
\midrule
IS 夏普均值 & $1.26$ & 样本内选出的组合平均夏普 \\
OOS 夏普均值 & $0.57$ & 同一组合在紧随其后的样本外平均夏普 \\
相对衰减 & $54.7\%$ & 超过一半的样本内表现不能兑现 \\
OOS 为正的折占比 & $44.4\%$ & 不到一半的样本外窗口录得正收益 \\
被选组合的 OOS 平均分位 & $0.66$ & 高于随机（$0.5$），选择过程有信息但弱 \\
IS$\to$OOS 秩相关 & $0.27$（$p=0.000$） & 统计显著但幅度低 \\
OOS 中位数组合夏普 & $-0.26$ & 随便挑一个组合的典型结果 \\
OOS 事后最优夏普 & $1.93$ & 事后诸葛的上界 \\
\bottomrule
\end{tabular}
\end{table}

\textbf{三点评读}。

第一，\textbf{选择过程确实携带信息，但远不足以覆盖衰减}。被选组合在样本外的平均分位 $0.66$、秩相关 $0.27$ 且统计显著（$p=0.000$），说明「样本内夏普高」与「样本外夏普高」之间不是噪声关系。但相对衰减 $54.7\%$、OOS 为正月数仅 $44.4\%$——选择带来的收益改善（$0.57$ 对中位数 $-0.26$，改善 $0.83$）被衰减吃掉了大半。

第二，\textbf{选中组合高度集中，且包含证据薄弱者}。$9$ 折中被选最多的三个是 \texttt{astk\_XS\_rev\_L120}（$3$ 次）、\texttt{astk\_XS\_rev\_L60}（$2$ 次）、\texttt{astk\_TS\_rev\_L250}（$2$ 次）。前两个属于论文表 $6$ 直接支持的口径，尚可；但 \texttt{astk\_TS\_rev\_L250} 属 A 股宽基指数时序口径——论文已明确标注该口径\emph{证据力弱}（$7$ 个指数相关性 $>0.85$，实质是单一 beta 择时）。它在两折中被选中，说明\emph{单纯按历史夏普选组合会选到证据薄弱的口径}。这正是多重检验流程要求登记搜索空间的原因。

第三，\textbf{这些数字不支持「按历史最优参数上线」}。论文结论置信度一章曾把「本文实证支持持有期决定归属」评为低置信度，理由是样本单一、未做多重检验、做空不可实现。walk-forward 给出了第四个理由，而且是最直接的：\emph{样本内最优参数在样本外平均衰减过半}。

\subsection{行业中性化口径：真实申万 vs 聚类代理}

前文规则二使用「收益相关结构做主成分 $+$ 聚类」得到的行业代理，因为当时未找到本地真实行业分类。本次从通达信自带文件解析出了真实分类：

\begin{itemize}[leftmargin=2.2em]
  \item \texttt{<TDX>/incon.dat}：行业码 $\to$ 行业名（含 \texttt{\#TDXNHY} 通达信行业、\texttt{\#TDXRSHY} 申万体系行业两段）；
  \item \texttt{<TDX>/T0002/hq\_cache/tdxyh.cfg}：证券代码 $\to$（通达信行业码, 申万体系行业码）。
\end{itemize}

解析结果：\textbf{申万体系一级 $30$ 个行业、覆盖 $5586$ 只证券}；二级 $128$ 个行业。实现见 \texttt{code/industry\_data.py}，无需联网。

\textbf{对照结果}（规则二纯多头版，$L=60$，其余参数不变）：

\begin{table}[htbp]
\centering\footnotesize
\caption{行业中性化口径对照}\label{tab:indcmp}
\begin{tabular}{@{}>{\raggedright\arraybackslash}p{3.6cm}>{\raggedright\arraybackslash}p{2.0cm}>{\raggedright\arraybackslash}p{2.0cm}>{\raggedright\arraybackslash}p{2.0cm}>{\raggedright\arraybackslash}p{2.4cm}@{}}
\toprule
\textbf{中性化口径} & \textbf{夏普} & \textbf{年化} & \textbf{换手} & \textbf{成本/毛利} \\
\midrule
真实申万一级（$30$ 行业） & $0.19$ & $1.69\%$ & $10.39$ & $32.22\%$ \\
聚类代理（$k=20$ 簇） & $0.23$ & $2.25\%$ & $9.02$ & $26.94\%$ \\
\bottomrule
\end{tabular}
\end{table}

\textbf{结论：聚类代理基本够用}。两者夏普相差 $0.04$，量级一致。真实申万版的换手更高（$10.39$ 对 $9.02$）、成本占比更高（$32.22\%$ 对 $26.94\%$），已经贴近铁律二的 $3$ 倍门槛（$33\%$）——这是一个需要留意的信号：\emph{用真实行业分类做中性化反而更接近成本门槛}。

但从严谨性看，真实分类应当作为默认：它可解释、可审计、不依赖聚类超参数（簇数 $k$ 与随机种子）。本文据此把规则二的默认口径从 \texttt{cluster} 改为 \texttt{sw}。

需要说明一项已知偏差：本地行业分类是\emph{静态快照}（当前时点的分类），不含行业变更历史。真实历史回测应当使用\emph{时点}行业分类；本文无法获得，因此该偏差保留。

\subsection{规则一完整窗口网格：工程增强改变了窗口排序}

前文只对 $L\in\{250,60,10\}$ 三档做了记录。本次跑完 $6$ 个窗口的完整网格，结果揭示了前文的一处判断错误。

\begin{table}[htbp]
\centering\footnotesize
\caption{规则一完整窗口网格：论文口径 vs 加入工程增强后}\label{tab:grid1}
\begin{tabular}{@{}>{\raggedright\arraybackslash}p{1.3cm}>{\raggedright\arraybackslash}p{1.9cm}>{\raggedright\arraybackslash}p{1.9cm}>{\raggedright\arraybackslash}p{2.1cm}>{\raggedright\arraybackslash}p{2.0cm}>{\raggedright\arraybackslash}p{2.0cm}@{}}
\toprule
\multirow{2}{*}{\textbf{$L$}} & \multicolumn{2}{c}{\textbf{夏普}} & \multicolumn{2}{c}{\textbf{换手}} & \multirow{2}{*}{\textbf{成本门槛}} \\
\cmidrule(lr){2-3}\cmidrule(lr){4-5}
 & 论文口径 & 规则一 & 论文口径 & 规则一 & \\
\midrule
$5$   & $+0.20$ & $+0.14$ & $52.41$ & $36.08$ & 不通过（$44.9\%$） \\
$10$  & $+0.32$ & $\mathbf{+0.44}$ & $35.14$ & $44.28$ & 通过（$18.1\%$） \\
$20$  & $+0.21$ & $\mathbf{+0.33}$ & $24.78$ & $33.68$ & 通过（$17.8\%$） \\
$60$  & $+0.21$ & $+0.19$ & $16.14$ & $22.98$ & 通过（$20.1\%$） \\
$120$ & $+0.09$ & $+0.05$ & $12.10$ & $18.74$ & 不通过（$44.5\%$） \\
$250$ & $+0.24$ & $\mathbf{-0.00}$ & $8.83$ & $12.74$ & \textbf{不通过（$108.7\%$）} \\
\bottomrule
\end{tabular}
\end{table}

\textbf{被推翻的判断}。前文把 $L=250$ 推荐为「低换手档」，依据是论文表 $5$ 的夏普 $0.24$、换手 $8.83$、成本占毛利 $10.2\%$。加入工程增强后，$L=250$ 的夏普降到 $-0.00$、成本占毛利升到 $108.7\%$——\textbf{该档失效}。

\textbf{成因}。工程增强项对长窗口的伤害最大，三个机制叠加：

\begin{enumerate}[leftmargin=2.2em]
  \item \textbf{入场门槛}：要求标准化信号强度 $|s|/(\hat\sigma\sqrt L)>1.0$。$L=250$ 时分母含 $\sqrt{250}$，门槛相对信号被大幅抬高，大量本可持仓的品种被弃权（实测平均持仓由论文的全池降到 $11.1$ 个品种）。
  \item \textbf{免交易带}：带宽 $\theta\propto V_{\text{target}}/\hat\sigma_i$ 与窗口无关，但长窗口的信号变化本来就小，宽带的相对抑制更强。
  \item \textbf{ATR 止损}：$L=250$ 是慢信号，$2.5$ 倍 $\text{ATR}_{20}$ 的止损会在信号兑现前把仓位砍掉（止损 $279$ 次）。
\end{enumerate}

三者共同把 $L=250$ 的毛收益压到接近零，而换手成本照付，于是成本占毛利比爆炸。

\textbf{修正后的推荐}。规则一的可用区间应从「$250/60/10$ 三档」改为「\textbf{$10$ 与 $20$ 两档}」：$L=10$ 夏普 $0.44$、$L=20$ 夏普 $0.33$，且二者成本占毛利均在 $18\%$ 以内、余量充足。$L=60$ 勉强可用（$0.19$），$L=120$ 与 $L=250$ 应放弃。

这一修正同时给出了一个方法论教训：\textbf{论文口径下的最优参数，在加入工程增强后可能不再最优，甚至失效}。工程增强必须与参数选择一起重做网格，不能沿用原论文的最优参数。

\subsection{规则三统计口径修正}

修正两处。

\subsubsection*{修正一：改用 Engle--Granger 专用临界值}

前文用标准 ADF 检验的 $5\%$ 临界值（约 $-2.86$）判定协整。这在方法上是错的：标准 ADF 适用于\emph{已知无协整关系}的单序列单位根检验；而协整检验的残差是\emph{估计出来的}，其分布左移，必须用更负的临界值。按 MacKinnon 响应面（$N=2$，含常数项）

\begin{equation}
\text{crit}_{5\%}(T)=-3.3377-\frac{5.967}{T}-\frac{8.98}{T^2},
\qquad T=60\ \Rightarrow\ \text{crit}=-3.4396 .
\end{equation}

\begin{table}[htbp]
\centering\footnotesize
\caption{S3 临界值口径修正的影响}\label{tab:egcrit}
\begin{tabular}{@{}>{\raggedright\arraybackslash}p{4.0cm}>{\raggedright\arraybackslash}p{2.6cm}>{\raggedright\arraybackslash}p{2.6cm}>{\raggedright\arraybackslash}p{2.6cm}@{}}
\toprule
\textbf{口径} & \textbf{通过条数} & \textbf{通过者 ADF 中位数} & \textbf{半衰期中位数} \\
\midrule
标准 ADF（$p<0.05$） & $75$ & $-3.25$ & $5.9$ 日 \\
EG 专用临界值（$-3.440$） & $26$ & $-3.68$ & $6.0$ 日 \\
\midrule
变化 & $-65.3\%$ & — & — \\
\bottomrule
\end{tabular}
\end{table}

\textbf{通过数从 $75$ 降到 $26$，减少 $65.3\%$}。也就是说，此前筛选出的「可交易对」中约三分之二是口径宽松造成的伪协整。前文把这一近似描述为「偏宽松、会略微高估通过率」——实测幅度是 $65\%$，远不止「略微」。这是本文对自身前文判断的一次修正。

\textbf{更严重的是样本量崩塌}。用 EG 严格口径重跑交易回测：

\begin{table}[htbp]
\centering\footnotesize
\caption{规则三在两种临界值口径下的结果}\label{tab:r3cmp}
\begin{tabular}{@{}>{\raggedright\arraybackslash}p{3.4cm}>{\raggedright\arraybackslash}p{2.2cm}>{\raggedright\arraybackslash}p{2.2cm}>{\raggedright\arraybackslash}p{2.2cm}>{\raggedright\arraybackslash}p{2.2cm}@{}}
\toprule
\textbf{口径} & \textbf{通过配对} & \textbf{成交笔数} & \textbf{夏普} & \textbf{年化} \\
\midrule
标准 ADF & $75$ & $21$ & $1.09$ & $4.31\%$ \\
EG 专用临界值 & $26$ & $\mathbf{7}$ & $0.48$ & $0.95\%$ \\
\bottomrule
\end{tabular}
\end{table}

\textbf{$7$ 笔交易在 $12$ 年样本上不足以支撑任何统计推断}。因此规则三在严格口径下的结论应当是：\emph{筛选流程可用，但样本量太小，策略无法被评估}。前文给出的「夏普 $1.09$」应作废——它建立在被证明过宽的口径上。

\subsubsection*{修正二：Kalman $q,r$ 参数敏感性}

\begin{table}[htbp]
\centering\footnotesize
\caption{Kalman $q,r$ 敏感性（EG 严格口径，$26$ 条配对）}\label{tab:kalsens}
\begin{tabular}{@{}>{\raggedright\arraybackslash}p{1.8cm}>{\raggedright\arraybackslash}p{1.6cm}>{\raggedright\arraybackslash}p{1.6cm}>{\raggedright\arraybackslash}p{1.7cm}>{\raggedright\arraybackslash}p{2.0cm}>{\raggedright\arraybackslash}p{1.8cm}@{}}
\toprule
\textbf{$q$} & \textbf{$r$} & \textbf{夏普} & \textbf{年化} & \textbf{成交笔数} & \textbf{平均持仓} \\
\midrule
$10^{-6}$ & $10^{-4}$ & $1.17$ & $1.10\%$ & $5$ & $14.0$ 日 \\
$10^{-6}$ & $10^{-2}$ & $0.70$ & $1.01\%$ & $4$ & $11.8$ 日 \\
$10^{-5}$ & $10^{-3}$ & $0.48$ & $0.95\%$ & $7$ & $12.4$ 日 \\
$10^{-4}$ & $10^{-3}$ & $1.06$ & $5.42\%$ & $22$ & $7.7$ 日 \\
$10^{-4}$ & $10^{-2}$ & $0.88$ & $4.69\%$ & $14$ & $4.9$ 日 \\
\bottomrule
\end{tabular}
\end{table}

（完整 $9$ 组见 \texttt{outputs/rule3\_kalman\_sensitivity.csv}。）

\textbf{夏普跨参数标准差 $0.20$，但成交笔数在 $3\sim22$ 之间波动（$7$ 倍）}。两点评读：

第一，夏普对 $q,r$ 相对不敏感（除非同时把两者都调小），说明结论的方向不受参数选择支配——这是好的一面。

第二，\textbf{交易频率对 $q,r$ 极度敏感}，而频率直接决定样本量（$3$ 笔 vs $22$ 笔）与成本。这意味着「规则三是否可评估」本身取决于一个此前未被注意的超参数。任何 $q,r$ 的选取都应当在报告结论时一并声明。

\subsection{四项验证对前文判断的修正汇总}

\begin{table}[htbp]
\centering\footnotesize
\caption{四项验证修正的判断}\label{tab:corrections}
\begin{tabular}{@{}>{\raggedright\arraybackslash}p{2.8cm}>{\raggedright\arraybackslash}p{5.0cm}>{\raggedright\arraybackslash}p{5.2cm}@{}}
\toprule
\textbf{验证} & \textbf{前文判断} & \textbf{修正后} \\
\midrule
walk-forward & 「可进入执行层对照」（流程第 6 步写得通过即放行） & IS$\to$OOS 衰减 $54.7\%$、仅 $44.4\%$ 折为正：不得按历史最优参数上线 \\
行业中性化 & 聚类代理为权宜之计，需接真实行业数据 & 真实申万（$0.19$）与聚类（$0.23$）量级一致，代理基本够用；但真实版更贴近成本门槛 \\
规则一窗口 & 推荐 $250/60/10$ 三档，$L=250$ 为低换手档 & $L=250$ 在工程增强下失效（$0.24\to-0.00$）；改为推荐 $L=10$ 与 $20$ \\
规则三临界值 & 「标准 ADF 偏宽松，会略微高估通过率」 & 高估幅度 $65.3\%$；严格口径下成交仅 $7$ 笔，策略无法评估 \\
\bottomrule
\end{tabular}
\end{table}

四项验证有一个共同的形态：\textbf{前文每一处「工程近似」或「暂用权宜」的地方，实跑后都被证明会实质改变结论}。这不是巧合，而是论文审计章那条规律在本文自身工作上的复现——「不同设置的差异使跨文献比较必须标注条件」，而同样的差异也出现在\emph{同一份工作内部}的不同版本之间。

据此，本文的定位收紧为\textbf{设计模板与验证流程}，而非可上线的策略；三条规则各自的失效事实与可用边界，见表~\ref{tab:corrections}~与第~\ref{sec:retune}~节的参数推荐总表。

\section{依据四项验证重标定参数}\label{sec:retune}

第~\ref{sec:verify}~节的四项验证各自推翻了前文的一处判断。本节据此重新标定规则一与规则二的参数，并给出必须同时声明的限定。

\subsection{重标定的三条判据}

参数推荐不能沿用论文口径的最优值（验证 3 已证明工程增强会改变最优参数），也不能只看全样本夏普（验证 1 已证明样本内表现不能外推）。本文采用三条判据，须\emph{同时}满足：

\begin{enumerate}[leftmargin=2.2em]
  \item \textbf{成本门槛}（铁律二）：毛利/成本 $>3$，即成本占毛利 $<33\%$。
  \item \textbf{参数曲面平滑}（验证 2 的 G7）：最优参数的邻域不得出现剧烈跳变。存在孤立峰者判为过拟合征兆，予以拒收。
  \item \textbf{样本外可兑现}（验证 1）：用规则自身的 walk-forward 检验参数选择是否外推。
\end{enumerate}

第 2 条是本次新增的判据，它在下面的扫描中起了决定性作用。

\subsection{规则一：窗口 $L$ 与入场门槛 $K$ 必须联合选择}

前文把入场门槛 $K=1.0$ 与窗口 $L$ 分别设定，这是错误的做法——验证 3 显示 $L=250$ 的失效与 $K=1.0$ 直接相关。因此本节做 $6\times4$ 的联合扫描（$L\in\{5,10,20,60,120,250\}$，$K\in\{0,0.5,1.0,1.5\}$）。

\begin{table}[htbp]
\centering\footnotesize
\caption{规则一 $L\times K$ 联合扫描（每格：夏普 / 年化\% / 成本占毛利\%）}\label{tab:r1LK}
\begin{tabular}{@{}>{\raggedright\arraybackslash}p{0.9cm}>{\raggedright\arraybackslash}p{2.5cm}>{\raggedright\arraybackslash}p{2.5cm}>{\raggedright\arraybackslash}p{2.5cm}>{\raggedright\arraybackslash}p{2.5cm}@{}}
\toprule
\textbf{$L$} & \textbf{$K=0$} & \textbf{$K=0.5$} & \textbf{$K=1.0$} & \textbf{$K=1.5$} \\
\midrule
$5$   & $0.29$ / $0.58$ / $3.6$  & $0.25$ / $0.90$ / $19.4$ & $0.14$ / $0.67$ / $44.9$ & $0.05$ / $-0.04$ / $67.6$ \\
$10$  & $\mathbf{0.56}$ / $0.90$ / $1.5$ & $0.52$ / $1.77$ / $7.2$ & $\mathbf{0.44}$ / $3.55$ / $18.1$ & $0.15$ / $1.10$ / $33.5$ \\
$20$  & $0.25$ / $0.56$ / $1.9$ & $0.12$ / $0.38$ / $19.4$ & $0.33$ / $2.58$ / $17.8$ & $0.25$ / $2.36$ / $19.4$ \\
$60$  & $0.22$ / $0.49$ / $1.4$ & $0.08$ / $0.25$ / $18.8$ & $0.19$ / $1.30$ / $20.1$ & $0.09$ / $0.18$ / $35.1$ \\
$120$ & $0.04$ / $0.08$ / $7.3$  & $0.06$ / $0.16$ / $16.4$ & $0.05$ / $0.08$ / $44.5$ & $-0.17$ / $-3.22$ / $-28.4$ \\
$250$ & $0.02$ / $0.02$ / $12.8$ & $-0.15$ / $-0.77$ / $-8.2$ & $-0.00$ / $-0.32$ / $108.7$ & $0.50$ / $5.63$ / $6.2$ \\
\bottomrule
\end{tabular}
\end{table}

这四行数据里最关键的信息不在数值高低，而在\textbf{曲面的形状}：

\begin{itemize}[leftmargin=2.2em]
  \item \textbf{$L=10$ 是唯一单调平滑的行}：$K$ 从 $0$ 升到 $1.5$，夏普依次 $0.56\to0.52\to0.44\to0.15$，单调递减、无跳变。这是稳健的参数结构——门槛越严，信号越少，夏普越低，符合预期。
  \item \textbf{$L=250$ 是非单调且有孤立峰的行}：$0.02\to-0.15\to-0.00\to\mathbf{0.50}$。$K=1.5$ 处的 $0.50$ 是孤立的好点，紧邻的 $K=1.0$ 却是 $-0.00$。更可疑的是$K=1.5$ 的换手（$22.64$）反而\emph{高于} $K=1.0$（$12.74$）——门槛更严却交易更频繁，说明信号在阈值附近反复穿越，这是典型的参数不稳定特征。\textbf{按判据 2，$L=250$ 全部拒收}。
  \item \textbf{$L=120$ 全档夏普 $\le0.06$}，直接排除。
  \item \textbf{$L=20$ 与 $L=60$ 非单调}（$K=0.5$ 处塌陷到 $0.12$ 与 $0.08$），稳健性弱于 $L=10$。
\end{itemize}

\textbf{修正后的规则一参数推荐}：首选 $L=10$、ENTRY\_K$=1.0$（夏普 $0.44$、年化 $3.55\%$、最大回撤 $-19.8\%$、成本占毛利 $18.1\%$），备选 $L=20$、ENTRY\_K$=1.0$（夏普 $0.33$、年化 $2.58\%$、回撤 $-16.9\%$）；$L=250$、$120$ 分别因曲面非单调与全档夏普过低排除，$L=5$ 在 $K\ge1.0$ 时不通过成本门槛。

为何首选 $K=1.0$ 而非夏普更高的 $K=0$（$0.56$）？因为 $K=0$ 意味着无入场门槛，全部品种按波动率倒数建仓，组合极度分散、波动极低，于是夏普虚高而年化仅 $0.90\%$。$K=1.0$ 的年化是其 $4$ 倍（$3.55\%$），夏普略低（$0.44$），成本占比 $18.1\%$ 仍在门槛内。这是收益与稳健性的折中。

\subsection{规则一：ATR 止损是低敏感参数}

前文把 ATR 止损倍数标为“工程判断”。本次做敏感性测试，结论是它可以不必精调。

\begin{table}[htbp]
\centering\footnotesize
\caption{规则一 ATR 止损倍数敏感性}\label{tab:r1atr}
\begin{tabular}{@{}>{\raggedright\arraybackslash}p{1.5cm}>{\raggedright\arraybackslash}p{1.9cm}>{\raggedright\arraybackslash}p{1.9cm}>{\raggedright\arraybackslash}p{2.1cm}>{\raggedright\arraybackslash}p{2.1cm}>{\raggedright\arraybackslash}p{2.1cm}@{}}
\toprule
\textbf{$L$} & \textbf{ATR 倍数} & \textbf{夏普} & \textbf{年化} & \textbf{最大回撤} & \textbf{止损次数} \\
\midrule
\multirow{4}{*}{$10$} & $1.5$ & $0.47$ & $3.63\%$ & $-17.83\%$ & $691$ \\
 & $2.5$ & $0.44$ & $3.55\%$ & $-19.82\%$ & $279$ \\
 & $4.0$ & $0.48$ & $4.01\%$ & $-16.98\%$ & $90$ \\
 & 关闭 & $0.45$ & $3.77\%$ & $-17.82\%$ & $0$ \\
\midrule
\multirow{4}{*}{$20$} & $1.5$ & $0.28$ & $2.03\%$ & $-17.21\%$ & $738$ \\
 & $2.5$ & $0.33$ & $2.58\%$ & $-16.91\%$ & $323$ \\
 & $4.0$ & $0.33$ & $2.64\%$ & $-17.50\%$ & $86$ \\
 & 关闭 & $0.32$ & $2.62\%$ & $-17.93\%$ & $0$ \\
\bottomrule
\end{tabular}
\end{table}

\textbf{$L=10$ 档的夏普跨度仅 $0.04$（$0.44\sim0.48$），$L=20$ 档仅 $0.05$（$0.28\sim0.33$）}。而止损次数从 $0$ 到 $738$ 相差两个数量级。这说明：

\begin{itemize}[leftmargin=2.2em]
  \item ATR 止损在本策略上\emph{既不创造也不破坏}收益，其主要作用是改变单品种的路径依赖与操作负荷；
  \item 该参数可以从“工程判断”降级为\textbf{低敏感参数}：取 $2.5$ 或直接关闭均可，无须精调；
  \item 若考虑操作成本（$691$ 次止损意味着大量小额交易与监控负担），取 $4.0$ 或关闭更省事，且夏普不劣。
\end{itemize}

这是四项验证中唯一“参数不确定性被消除”的结论，值得单独记下：\emph{并非所有未经验证的工程参数都重要，实测能把无关紧要的挑出来。}

\subsection{规则二：$L=250$ 的风险特征优于 $L=60$}

前文规则二只有 $L=60$ 一个数据点（来自论文口径），无法给出参数推荐。本次在真实申万行业口径下做 $L$ 扫描。

\begin{table}[htbp]
\centering\footnotesize
\caption{规则二窗口 $L$ 扫描（真实申万一级口径，纯多头）}\label{tab:r2L}
\begin{tabular}{@{}>{\raggedright\arraybackslash}p{1.3cm}>{\raggedright\arraybackslash}p{1.9cm}>{\raggedright\arraybackslash}p{1.9cm}>{\raggedright\arraybackslash}p{2.1cm}>{\raggedright\arraybackslash}p{1.7cm}>{\raggedright\arraybackslash}p{2.0cm}>{\raggedright\arraybackslash}p{1.8cm}@{}}
\toprule
\textbf{$L$} & \textbf{多头夏普} & \textbf{年化} & \textbf{最大回撤} & \textbf{换手} & \textbf{成本占毛利} & \textbf{门槛} \\
\midrule
$20$  & $0.09$ & $0.09\%$ & $-67.05\%$ & $12.83$ & $53.94\%$ & 不通过 \\
$60$  & $0.19$ & $1.69\%$ & $-61.54\%$ & $10.39$ & $32.22\%$ & 通过（贴线） \\
$120$ & $-0.15$ & $-2.93\%$ & $-64.94\%$ & $7.77$ & $-99.48\%$ & 收益为负 \\
$\mathbf{250}$ & $\mathbf{0.18}$ & $1.57\%$ & $\mathbf{-45.24\%}$ & $\mathbf{7.10}$ & $\mathbf{25.65\%}$ & \textbf{通过（余量足）} \\
\bottomrule
\end{tabular}
\end{table}

\textbf{$L=60$ 与 $L=250$ 的夏普几乎相同（$0.19$ 对 $0.18$），但 $L=250$ 在三个风险维度上全面更优}：

\begin{itemize}[leftmargin=2.2em]
  \item \textbf{最大回撤浅 $16$ 个百分点}（$-45.24\%$ 对 $-61.54\%$）；
  \item \textbf{换手低 $32\%$}（$7.10$ 对 $10.39$）；
  \item \textbf{成本占毛利低 $6.6$ 个百分点}（$25.65\%$ 对 $32.22\%$），且 $L=60$ 的 $32.22\%$ 已贴上 $33\%$ 门槛线，成本一旦上升即失效。
\end{itemize}

因此\textbf{规则二的推荐改为：$L=250$ 首选，$L=60$ 备选}。这与论文口径（A 股截面反转 $L=60$ 最优、$L=250$ 明显衰减）相反，是工程增强再次改变参数排序的案例——持仓上限、止损与回撤控制对长窗口更友好，因为长窗口的信号更稳、更耐受这些干预。

同时要注意，$L=20$ 不通过成本门槛（$53.94\%$），$L=120$ 收益为负，两者都应排除。可用区间因此是 $\{60, 250\}$，比论文口径的 $\{20,\dots,250\}$ 窄得多。

\textbf{一条必须保留的限定}：即便在最优的 $L=250$ 上，纯多头版本的对冲后夏普仍为 $-0.04$（第~\ref{sec:ashare}~节已指出其跑输被对冲的指数）。\emph{参数重标定改善的是相对排序，不改变该版本在样本期上失效的事实。}

\subsection{规则自身参数的样本外检验}

上面的推荐依据全样本夏普与曲面形状。更直接的检验是：\emph{若在样本内选参数，样本外能否兑现}。为此对规则自身的参数（规则一的 $L$、规则二的 $L$）做 walk-forward，做法是对每个候选参数跑一次全样本回测、拿到逐日净收益序列，再按滚动折切片（IS $756$ 日选参数、OOS $252$ 日检验）。

\begin{table}[htbp]
\centering\footnotesize
\caption{规则自身参数的 walk-forward 汇总（$9$ 折）}\label{tab:rulewf}
\begin{tabular}{@{}>{\raggedright\arraybackslash}p{3.0cm}>{\raggedright\arraybackslash}p{1.8cm}>{\raggedright\arraybackslash}p{1.8cm}>{\raggedright\arraybackslash}p{2.0cm}>{\raggedright\arraybackslash}p{1.8cm}>{\raggedright\arraybackslash}p{2.6cm}@{}}
\toprule
\textbf{规则} & \textbf{IS 夏普} & \textbf{OOS 夏普} & \textbf{OOS 为正折} & \textbf{选择增益} & \textbf{选中参数分布} \\
\midrule
规则一 & $0.74$ & $-0.03$ & $7/9$ & $+0.05$ & $\{L{=}10:6,\ L{=}20:2,\ L{=}5:1\}$ \\
规则二 & $0.20$ & $-0.18$ & $4/9$ & $+0.12$ & $\{L{=}250:5,\ L{=}120:3,\ L{=}60:1\}$ \\
\bottomrule
\end{tabular}
\end{table}

（“选择增益” = 所选参数的 OOS 夏普 $-$ 该折全部候选参数的 OOS 夏普中位数。完整逐折数据见 \texttt{outputs/rule\_wf\_rule1.csv} 与 \texttt{rule\_wf\_rule2.csv}。）

\textbf{两条规则的参数选择都稳定，但都不兑现收益}。

\textbf{规则一}：$9$ 折中有 $6$ 折选中 $L=10$、$2$ 折选中 $L=20$、$1$ 折选中 $L=5$——\emph{从未选中 $60/120/250$}。这与本节的推荐（$L=10$ 首选、$L=20$ 备选）完全一致，说明参数选择过程稳定。但 IS 均值 $0.74$ → OOS 均值 $-0.03$（相对衰减 $104.5\%$），选择增益仅 $+0.05$。其中第 $1$ 折（选中 $L=5$）的 OOS 为 $-2.12$，是极端值；剔除后 OOS 均值回升到 $+0.23$——\emph{结果对单一折高度敏感，样本量（$9$ 折）不足以支撑稳定结论}。

\textbf{规则二}：$9$ 折中有 $5$ 折选中 $L=250$、$3$ 折选中 $L=120$、$1$ 折选中 $L=60$。对 $L=250$ 的偏好与本节的推荐一致（这构成了对“$L=250$ 优于 $L=60$”的独立支持）。但 IS 均值仅 $0.20$ → OOS 均值 $-0.18$（相对衰减 $194.4\%$），$9$ 折中仅 $4$ 折为正，选择增益 $+0.12$。

\subsection{修正后的参数推荐总表}

\begin{table}[htbp]
\centering\footnotesize
\caption{参数推荐：原版与重标定后对照}\label{tab:retunecmp}
\begin{tabular}{@{}>{\raggedright\arraybackslash}p{2.2cm}>{\raggedright\arraybackslash}p{4.3cm}>{\raggedright\arraybackslash}p{5.8cm}@{}}
\toprule
\textbf{参数} & \textbf{前文推荐} & \textbf{重标定后} \\
\midrule
\multicolumn{3}{@{}l}{\textbf{规则一（期货时序趋势）}} \\
窗口 $L$ & $250$ / $60$ / $10$ 三档 & $\mathbf{10}$（首选）、$\mathbf{20}$（备选）；排除 $250$、$120$、$5$ \\
入场门槛 ENTRY\_K & $1.0$（独立设定） & $\mathbf{1.0}$，但必须与 $L$ \emph{联合}选择（$L=250$ 的失效源于 $K=1.0$） \\
调仓频率 & $5$ 交易日 & 不变（$5$ 交易日） \\
波动率目标 & $10\%$ & 不变 \\
ATR 止损倍数 & $2.5$（工程判断） & $2.5$ 或关闭均可（\textbf{低敏感}，夏普跨度 $0.04$） \\
恐慌阈值/拥挤阈值 & $-5\%$ / $90$ 分位 & 不变（未做敏感性，仍属未验证） \\
\midrule
\multicolumn{3}{@{}l}{\textbf{规则二（A 股截面反转）}} \\
窗口 $L$ & $60$（备选 $120$） & $\mathbf{250}$（首选）、$\mathbf{60}$（备选）；排除 $20$、$120$ \\
行业中性化口径 & 聚类代理（权宜） & $\mathbf{sw}$ 真实申万一级（与聚类代理量级一致，但可解释可审计） \\
持仓上限 & $50$ 只 & 不变（$50$ 只） \\
单票止损 & $15\%$（高波动 $10\%$） & 不变 \\
现金冗余 & $30\%$ & 不变 \\
\bottomrule
\end{tabular}
\end{table}

\textbf{代码中的落地}。上述推荐已写入回测框架，不再是文档与代码两张皮：\texttt{code/rule1\_trend.py} 与 \texttt{code/rule2\_ashare.py} 各自定义了 \texttt{RECOMMENDED}（首选/备选两档，含实测指标注释）与 \texttt{EXCLUDED}（被排除窗口及原因）两个常量，类属性 \texttt{Rule1Config.L}$=10$、\texttt{Rule2Config.L}$=250$ 亦同步为首选档，并提供 \texttt{--preset} 开关：

\begin{quote}\ttfamily\small
python rule1\_trend.py --preset 首选 \hfill \# L=10, ENTRY\_K=1.0\\
python rule1\_trend.py --preset 备选 \hfill \# L=20, ENTRY\_K=1.0\\
python rule2\_ashare.py --preset 首选 \hfill \# L=250, 申万一级中性化\\
python rule1\_trend.py --grid \hfill \# 完整 6 窗口网格（复现第 11.3 节）
\end{quote}

两个首选档均已实跑复核，输出与本节表格逐项一致（规则一：夏普 $0.44$ / 年化 $3.55\%$ / 回撤 $-19.82\%$ / 成本占毛利 $18.09\%$；规则二：多头夏普 $0.18$ / 年化 $1.57\%$ / 回撤 $-45.24\%$ / 成本占毛利 $25.65\%$）。规则三的 \texttt{USE\_EG\_CRIT} 亦默认为 \texttt{True}（EG 专用临界值），与第~\ref{sec:verify}~节的修正一致。

\subsection{必须同时声明的两条限定}

重标定给出了更稳健的参数，但有两件事不能因此被掩盖。

\textbf{其一，重标定改善的是相对排序，不是绝对水平}。两条规则的首选参数在自身 walk-forward 中的样本外夏普分别为 $-0.03$ 与 $-0.18$（表~\ref{tab:rulewf}）——\textbf{本文推荐的参数在样本外均不兑现正收益}。推荐的依据是“参数曲面平滑、选择过程稳定”，不是“能赚钱”。

\textbf{其二，可推荐性建立在很薄的样本上}。规则一的 walk-forward 只有 $9$ 折且对单一折高度敏感；规则二的 IS 夏普本身只有 $0.20$，远低于论文口径的 $0.99$（差异来自流动性过滤与纯多头化）。

因此本次重标定的正确读法是：\textbf{它把参数从“论文口径的最优值”修正为“在本文工程口径下更稳健的值”，但并未使规则达到可上线的标准}。第~\ref{sec:verify}~节给出的定位——设计模板与验证流程——在重标定后依然成立。

\section*{参考文献}
\addcontentsline{toc}{section}{参考文献}
\begin{thebibliography}{99}
\small
\bibitem{shleifervishny1997} Shleifer, A. \& Vishny, R. W. (1997). \emph{The Limits of Arbitrage}. Journal of Finance 52(1): 35--55. \href{https://doi.org/10.1111/j.1540-6261.1997.tb03807.x}{https://doi.org/10.1111/j.1540-6261.1997.tb03807.x}.
\bibitem{khandanilo2011} Khandani, A. E. \& Lo, A. W. (2011). \emph{What Happened to the Quants in August 2007? Evidence from Factors and Transactions Data}. Journal of Financial Markets 14(1): 1--46. \href{https://doi.org/10.1016/j.finmar.2010.07.005}{https://doi.org/10.1016/j.finmar.2010.07.005}.
\bibitem{brunnermeier2009} Brunnermeier, M. K. \& Pedersen, L. H. (2009). \emph{Market Liquidity and Funding Liquidity}. Review of Financial Studies 22(6): 2201--2238. \href{https://doi.org/10.1093/rfs/hhn098}{https://doi.org/10.1093/rfs/hhn098}.
\bibitem{mop2012} Moskowitz, T. J., Ooi, Y. H. \& Pedersen, L. H. (2012). \emph{Time Series Momentum}. Journal of Financial Economics 104(2): 228--250. \href{https://doi.org/10.1016/j.jfineco.2011.11.003}{https://doi.org/10.1016/j.jfineco.2011.11.003}.
\bibitem{hop2017} Hurst, B., Ooi, Y. H. \& Pedersen, L. H. (2017). \emph{A Century of Evidence on Trend-Following Investing}. Journal of Portfolio Management 44(1): 15--29. \href{https://doi.org/10.3905/jpm.2017.44.1.015}{https://doi.org/10.3905/jpm.2017.44.1.015}.
\bibitem{barroso2015} Barroso, P. \& Santa-Clara, P. (2015). \emph{Momentum Has Its Moments}. Journal of Financial Economics 116(1): 111--120. \href{https://doi.org/10.1016/j.jfineco.2014.11.010}{https://doi.org/10.1016/j.jfineco.2014.11.010}.
\bibitem{harvey2018} Harvey, C. R., Hoyle, E., Korgaonkar, R., Rattray, S., Sargaison, M. \& van Hemert, O. (2018). \emph{The Impact of Volatility Targeting}. Journal of Portfolio Management 45(1): 14--33. \href{https://doi.org/10.3905/jpm.2018.45.1.014}{https://doi.org/10.3905/jpm.2018.45.1.014}.
\bibitem{moreiramuir2017} Moreira, A. \& Muir, T. (2017). \emph{Volatility-Managed Portfolios}. Journal of Finance 72(4): 1611--1644. \href{https://doi.org/10.1111/jofi.12513}{https://doi.org/10.1111/jofi.12513}.
\bibitem{kim2016} Kim, A. Y., Tse, Y. \& Wald, J. K. (2016). \emph{Time Series Momentum and Volatility Scaling}. Journal of Financial Markets 30: 103--124. \href{https://doi.org/10.1016/j.finmar.2016.05.003}{https://doi.org/10.1016/j.finmar.2016.05.003}.
\bibitem{huang2020} Huang, D., Li, J., Wang, L. \& Zhou, G. (2020). \emph{Time Series Momentum: Is It There?}. Journal of Financial Economics 135(3): 774--794. \href{https://doi.org/10.1016/j.jfineco.2019.08.004}{https://doi.org/10.1016/j.jfineco.2019.08.004}.
\bibitem{danielmoskowitz2016} Daniel, K. \& Moskowitz, T. J. (2016). \emph{Momentum Crashes}. Journal of Financial Economics 122(2): 221--247. \href{https://doi.org/10.1016/j.jfineco.2015.12.002}{https://doi.org/10.1016/j.jfineco.2015.12.002}.
\bibitem{stein2009} Stein, J. C. (2009). \emph{Presidential Address: Sophisticated Investors and Market Efficiency}. Journal of Finance 64(4): 1517--1548. \href{https://doi.org/10.1111/j.1540-6261.2009.01472.x}{https://doi.org/10.1111/j.1540-6261.2009.01472.x}.
\bibitem{loupolk2022} Lou, D. \& Polk, C. (2022). \emph{Comomentum: Inferring Arbitrage Activity from Return Correlations}. Review of Financial Studies 35(2): 735--776. \href{https://doi.org/10.1093/rfs/hhab117}{https://doi.org/10.1093/rfs/hhab117}.
\bibitem{baltas2019} Baltas, N. (2019). \emph{The Impact of Crowding in Alternative Risk Premia Investing}. Financial Analysts Journal 75(3): 89--104. \href{https://doi.org/10.1080/0015198x.2019.1600955}{https://doi.org/10.1080/0015198x.2019.1600955}.
\bibitem{debondt1985} De Bondt, W. F. M. \& Thaler, R. (1985). \emph{Does the Stock Market Overreact?}. Journal of Finance 40(3): 793--805. \href{https://doi.org/10.1111/j.1540-6261.1985.tb05004.x}{https://doi.org/10.1111/j.1540-6261.1985.tb05004.x}.
\bibitem{lakonishok1994} Lakonishok, J., Shleifer, A. \& Vishny, R. W. (1994). \emph{Contrarian Investment, Extrapolation, and Risk}. Journal of Finance 49(5): 1541--1578. \href{https://doi.org/10.1111/j.1540-6261.1994.tb04772.x}{https://doi.org/10.1111/j.1540-6261.1994.tb04772.x}.
\bibitem{moskowitzgrinblatt1999} Moskowitz, T. J. \& Grinblatt, M. (1999). \emph{Do Industries Explain Momentum?}. Journal of Finance 54(4): 1249--1290. \href{https://doi.org/10.1111/0022-1082.00146}{https://doi.org/10.1111/0022-1082.00146}.
\bibitem{leungli2015} Leung, T. \& Li, X. (2015). \emph{Optimal Mean Reversion Trading with Transaction Costs and Stop-Loss Exit}. International Journal of Theoretical and Applied Finance 18(3): 1550020. \href{https://doi.org/10.1142/S021902491550020X}{https://doi.org/10.1142/S021902491550020X}.
\bibitem{vidyamurthy2004} Vidyamurthy, G. (2004). \emph{Pairs Trading: Quantitative Methods and Analysis}. John Wiley \& Sons (ISBN 978-0-471-46067-1). \href{https://www.wiley.com/}{https://www.wiley.com/}.
\bibitem{krauss2017} Krauss, C. (2017). \emph{Statistical Arbitrage Pairs Trading Strategies: Review and Outlook}. Journal of Economic Surveys 31(2): 513--545. \href{https://doi.org/10.1111/joes.12153}{https://doi.org/10.1111/joes.12153}.
\bibitem{gatev2006} Gatev, E., Goetzmann, W. N. \& Rouwenhorst, K. G. (2006). \emph{Pairs Trading: Performance of a Relative-Value Arbitrage Rule}. Review of Financial Studies 19(3): 797--827. \href{https://doi.org/10.1093/rfs/hhj020}{https://doi.org/10.1093/rfs/hhj020}.
\bibitem{englegranger1987} Engle, R. F. \& Granger, C. W. J. (1987). \emph{Co-Integration and Error Correction: Representation, Estimation, and Testing}. Econometrica 55(2): 251--276. \href{https://doi.org/10.2307/1913236}{https://doi.org/10.2307/1913236}.
\bibitem{dickeyfuller1979} Dickey, D. A. \& Fuller, W. A. (1979). \emph{Distribution of the Estimators for Autoregressive Time Series With a Unit Root}. JASA 74(366): 427--431. \href{https://doi.org/10.1080/01621459.1979.10482531}{https://doi.org/10.1080/01621459.1979.10482531}.
\bibitem{kendall1953} Kendall, M. G. (1953). \emph{The Analysis of Economic Time-Series---Part I: Prices}. J. Royal Statistical Society A 116(1): 11--34. \href{https://doi.org/10.2307/2980947}{https://doi.org/10.2307/2980947}.
\bibitem{elliott2005} Elliott, R. J., van der Hoek, J. \& Malcolm, W. P. (2005). \emph{Pairs Trading}. Quantitative Finance 5(3): 271--276. \href{https://doi.org/10.1080/14697680500149370}{https://doi.org/10.1080/14697680500149370}.
\bibitem{martin2011} Martin, R. \& Schöneborn, T. (2011). \emph{Mean Reversion Pays, but Costs}. Risk 24(2): 84--89；arXiv:\allowbreak 1103.4934. \href{https://arxiv.org/abs/1103.4934}{https://arxiv.org/abs/1103.4934}.
\bibitem{avellaneda2010} Avellaneda, M. \& Lee, J.-H. (2010). \emph{Statistical Arbitrage in the US Equities Market}. Quantitative Finance 10(7): 761--782. \href{https://doi.org/10.1080/14697680903124632}{https://doi.org/10.1080/14697680903124632}.
\bibitem{almgrenchriss2000} Almgren, R. \& Chriss, N. (2000). \emph{Optimal Execution of Portfolio Transactions}. The Journal of Risk 3(2): 5--39. \href{https://doi.org/10.21314/jor.2001.041}{https://doi.org/10.21314/jor.2001.041}.
\bibitem{perold1988} Perold, A. F. (1988). \emph{The Implementation Shortfall: Paper versus Reality}. Journal of Portfolio Management 14(3): 4--9. \href{https://doi.org/10.3905/jpm.1988.409150}{https://doi.org/10.3905/jpm.1988.409150}.
\bibitem{white2000} White, H. (2000). \emph{A Reality Check for Data Snooping}. Econometrica 68(5): 1097--1126. \href{https://doi.org/10.1111/1468-0262.00152}{https://doi.org/10.1111/1468-0262.00152}.
\bibitem{hansen2005} Hansen, P. R. (2005). \emph{A Test for Superior Predictive Ability}. Journal of Business \& Economic Statistics 23(4): 365--380. \href{https://doi.org/10.1198/073500105000000063}{https://doi.org/10.1198/073500105000000063}.
\bibitem{romanowolf2005} Romano, J. P. \& Wolf, M. (2005). \emph{Stepwise Multiple Testing as Formalized Data Snooping}. Econometrica 73(4): 1237--1282. \href{https://doi.org/10.1111/j.1468-0262.2005.00615.x}{https://doi.org/10.1111/j.1468-0262.2005.00615.x}.
\bibitem{bailey2014} Bailey, D. H. \& López de Prado, M. (2014). \emph{The Deflated Sharpe Ratio}. Journal of Portfolio Management 40(5): 94--107. \href{https://doi.org/10.3905/jpm.2014.40.5.094}{https://doi.org/10.3905/jpm.2014.40.5.094}.
\bibitem{bailey2016} Bailey, D. H., Borwein, J. M., López de Prado, M. \& Zhu, Q. J. (2016). \emph{The Probability of Backtest Overfitting}. The Journal of Computational Finance 20(4): 39--69. \href{https://doi.org/10.21314/jcf.2016.322}{https://doi.org/10.21314/jcf.2016.322}.
\end{thebibliography}


\end{document}